% Options for packages loaded elsewhere
\PassOptionsToPackage{unicode}{hyperref}
\PassOptionsToPackage{hyphens}{url}
\documentclass[
]{article}
\usepackage{xcolor}
\usepackage[margin=1in]{geometry}
\usepackage{amsmath,amssymb}
\setcounter{secnumdepth}{-\maxdimen} % remove section numbering
\usepackage{iftex}
\ifPDFTeX
  \usepackage[T1]{fontenc}
  \usepackage[utf8]{inputenc}
  \usepackage{textcomp} % provide euro and other symbols
\else % if luatex or xetex
  \usepackage{unicode-math} % this also loads fontspec
  \defaultfontfeatures{Scale=MatchLowercase}
  \defaultfontfeatures[\rmfamily]{Ligatures=TeX,Scale=1}
\fi
\usepackage{lmodern}
\ifPDFTeX\else
  % xetex/luatex font selection
\fi
% Use upquote if available, for straight quotes in verbatim environments
\IfFileExists{upquote.sty}{\usepackage{upquote}}{}
\IfFileExists{microtype.sty}{% use microtype if available
  \usepackage[]{microtype}
  \UseMicrotypeSet[protrusion]{basicmath} % disable protrusion for tt fonts
}{}
\makeatletter
\@ifundefined{KOMAClassName}{% if non-KOMA class
  \IfFileExists{parskip.sty}{%
    \usepackage{parskip}
  }{% else
    \setlength{\parindent}{0pt}
    \setlength{\parskip}{6pt plus 2pt minus 1pt}}
}{% if KOMA class
  \KOMAoptions{parskip=half}}
\makeatother
\usepackage{color}
\usepackage{fancyvrb}
\newcommand{\VerbBar}{|}
\newcommand{\VERB}{\Verb[commandchars=\\\{\}]}
\DefineVerbatimEnvironment{Highlighting}{Verbatim}{commandchars=\\\{\}}
% Add ',fontsize=\small' for more characters per line
\usepackage{framed}
\definecolor{shadecolor}{RGB}{248,248,248}
\newenvironment{Shaded}{\begin{snugshade}}{\end{snugshade}}
\newcommand{\AlertTok}[1]{\textcolor[rgb]{0.94,0.16,0.16}{#1}}
\newcommand{\AnnotationTok}[1]{\textcolor[rgb]{0.56,0.35,0.01}{\textbf{\textit{#1}}}}
\newcommand{\AttributeTok}[1]{\textcolor[rgb]{0.13,0.29,0.53}{#1}}
\newcommand{\BaseNTok}[1]{\textcolor[rgb]{0.00,0.00,0.81}{#1}}
\newcommand{\BuiltInTok}[1]{#1}
\newcommand{\CharTok}[1]{\textcolor[rgb]{0.31,0.60,0.02}{#1}}
\newcommand{\CommentTok}[1]{\textcolor[rgb]{0.56,0.35,0.01}{\textit{#1}}}
\newcommand{\CommentVarTok}[1]{\textcolor[rgb]{0.56,0.35,0.01}{\textbf{\textit{#1}}}}
\newcommand{\ConstantTok}[1]{\textcolor[rgb]{0.56,0.35,0.01}{#1}}
\newcommand{\ControlFlowTok}[1]{\textcolor[rgb]{0.13,0.29,0.53}{\textbf{#1}}}
\newcommand{\DataTypeTok}[1]{\textcolor[rgb]{0.13,0.29,0.53}{#1}}
\newcommand{\DecValTok}[1]{\textcolor[rgb]{0.00,0.00,0.81}{#1}}
\newcommand{\DocumentationTok}[1]{\textcolor[rgb]{0.56,0.35,0.01}{\textbf{\textit{#1}}}}
\newcommand{\ErrorTok}[1]{\textcolor[rgb]{0.64,0.00,0.00}{\textbf{#1}}}
\newcommand{\ExtensionTok}[1]{#1}
\newcommand{\FloatTok}[1]{\textcolor[rgb]{0.00,0.00,0.81}{#1}}
\newcommand{\FunctionTok}[1]{\textcolor[rgb]{0.13,0.29,0.53}{\textbf{#1}}}
\newcommand{\ImportTok}[1]{#1}
\newcommand{\InformationTok}[1]{\textcolor[rgb]{0.56,0.35,0.01}{\textbf{\textit{#1}}}}
\newcommand{\KeywordTok}[1]{\textcolor[rgb]{0.13,0.29,0.53}{\textbf{#1}}}
\newcommand{\NormalTok}[1]{#1}
\newcommand{\OperatorTok}[1]{\textcolor[rgb]{0.81,0.36,0.00}{\textbf{#1}}}
\newcommand{\OtherTok}[1]{\textcolor[rgb]{0.56,0.35,0.01}{#1}}
\newcommand{\PreprocessorTok}[1]{\textcolor[rgb]{0.56,0.35,0.01}{\textit{#1}}}
\newcommand{\RegionMarkerTok}[1]{#1}
\newcommand{\SpecialCharTok}[1]{\textcolor[rgb]{0.81,0.36,0.00}{\textbf{#1}}}
\newcommand{\SpecialStringTok}[1]{\textcolor[rgb]{0.31,0.60,0.02}{#1}}
\newcommand{\StringTok}[1]{\textcolor[rgb]{0.31,0.60,0.02}{#1}}
\newcommand{\VariableTok}[1]{\textcolor[rgb]{0.00,0.00,0.00}{#1}}
\newcommand{\VerbatimStringTok}[1]{\textcolor[rgb]{0.31,0.60,0.02}{#1}}
\newcommand{\WarningTok}[1]{\textcolor[rgb]{0.56,0.35,0.01}{\textbf{\textit{#1}}}}
\usepackage{graphicx}
\makeatletter
\newsavebox\pandoc@box
\newcommand*\pandocbounded[1]{% scales image to fit in text height/width
  \sbox\pandoc@box{#1}%
  \Gscale@div\@tempa{\textheight}{\dimexpr\ht\pandoc@box+\dp\pandoc@box\relax}%
  \Gscale@div\@tempb{\linewidth}{\wd\pandoc@box}%
  \ifdim\@tempb\p@<\@tempa\p@\let\@tempa\@tempb\fi% select the smaller of both
  \ifdim\@tempa\p@<\p@\scalebox{\@tempa}{\usebox\pandoc@box}%
  \else\usebox{\pandoc@box}%
  \fi%
}
% Set default figure placement to htbp
\def\fps@figure{htbp}
\makeatother
\setlength{\emergencystretch}{3em} % prevent overfull lines
\providecommand{\tightlist}{%
  \setlength{\itemsep}{0pt}\setlength{\parskip}{0pt}}
\usepackage{bookmark}
\IfFileExists{xurl.sty}{\usepackage{xurl}}{} % add URL line breaks if available
\urlstyle{same}
\hypersetup{
  pdftitle={Mobilité \& Économie comportementale},
  pdfauthor={Kayliah Avorebang Kombila \textbar{} Yasmine Aouaj \textbar{} Lily Porre \textbar{} Melissa Mariano Albuquerque},
  hidelinks,
  pdfcreator={LaTeX via pandoc}}

\title{Mobilité \& Économie comportementale}
\author{Kayliah Avorebang Kombila \textbar{} Yasmine Aouaj \textbar{}
Lily Porre \textbar{} Melissa Mariano Albuquerque}
\date{2026-05-16}

\begin{document}
\maketitle

\section{PROJET : MOBILITÉ \& ÉCONOMIE COMPORTEMENTALE (FRANCE
2019-2026)}\label{projet-mobilituxe9-uxe9conomie-comportementale-france-2019-2026}

\subsection{Prix de l'essence \& habitudes de
déplacement}\label{prix-de-lessence-habitudes-de-duxe9placement}

\begin{Shaded}
\begin{Highlighting}[]
\NormalTok{packages }\OtherTok{\textless{}{-}} \FunctionTok{c}\NormalTok{(}
  \StringTok{"readxl"}\NormalTok{,}
  \StringTok{"dplyr"}\NormalTok{,}
  \StringTok{"tidyr"}\NormalTok{,}
  \StringTok{"purrr"}\NormalTok{,}
  \StringTok{"broom"}\NormalTok{,}
  \StringTok{"ggplot2"}\NormalTok{,}
  \StringTok{"ggpmisc"}\NormalTok{,}
  \StringTok{"patchwork"}\NormalTok{,}
  \StringTok{"scales"}
\NormalTok{)}

\NormalTok{installed  }\OtherTok{\textless{}{-}} \FunctionTok{rownames}\NormalTok{(}\FunctionTok{installed.packages}\NormalTok{())}
\NormalTok{to\_install }\OtherTok{\textless{}{-}}\NormalTok{ packages[}\SpecialCharTok{!}\NormalTok{packages }\SpecialCharTok{\%in\%}\NormalTok{ installed]}
\ControlFlowTok{if}\NormalTok{ (}\FunctionTok{length}\NormalTok{(to\_install) }\SpecialCharTok{\textgreater{}} \DecValTok{0}\NormalTok{) \{}
  \FunctionTok{install.packages}\NormalTok{(to\_install, }\AttributeTok{repos =} \StringTok{"https://cran.r{-}project.org"}\NormalTok{)}
\NormalTok{\}}

\FunctionTok{invisible}\NormalTok{(}\FunctionTok{lapply}\NormalTok{(packages, library, }\AttributeTok{character.only =} \ConstantTok{TRUE}\NormalTok{))}
\end{Highlighting}
\end{Shaded}

\subsubsection{jeu de donées}\label{jeu-de-donuxe9es}

\begin{Shaded}
\begin{Highlighting}[]
\NormalTok{df\_raw }\OtherTok{\textless{}{-}} \FunctionTok{read.csv2}\NormalTok{(}
  \StringTok{"dataset\_france.csv"}\NormalTok{,}
  \AttributeTok{stringsAsFactors =} \ConstantTok{TRUE}
\NormalTok{)}
\end{Highlighting}
\end{Shaded}

\begin{Shaded}
\begin{Highlighting}[]
\CommentTok{\# Dimensions du dataset}
\FunctionTok{dim}\NormalTok{(df\_raw)}
\end{Highlighting}
\end{Shaded}

\begin{verbatim}
## [1] 274  22
\end{verbatim}

L'importation du fichier donne un ensemble de 274 observations et 22
variables.

\begin{Shaded}
\begin{Highlighting}[]
\CommentTok{\# Structure des variables}
\FunctionTok{str}\NormalTok{(df\_raw)}
\end{Highlighting}
\end{Shaded}

\begin{verbatim}
## 'data.frame':    274 obs. of  22 variables:
##  $ Date              : Factor w/ 93 levels "","2019-01","2019-02",..: 2 2 2 3 3 3 4 4 4 5 ...
##  $ Zone              : Factor w/ 8 levels "","0,41","0,50",..: 8 6 7 8 6 7 8 6 7 8 ...
##  $ Prix_SP95         : Factor w/ 44 levels "","1,250 €","1,300 €",..: 9 9 9 10 10 10 11 11 11 15 ...
##  $ Prix.Gazole       : Factor w/ 50 levels "","0,38","0,47",..: 14 14 14 16 16 16 18 18 18 21 ...
##  $ Trafic.Routier    : Factor w/ 260 levels "","10001","10002",..: 96 22 198 82 18 188 112 42 212 121 ...
##  $ Voyageurs.TC      : Factor w/ 231 levels "","10073","10172",..: 200 55 129 207 61 133 222 76 142 3 ...
##  $ Index.SP95        : Factor w/ 43 levels "","0,62","0,64",..: 39 39 39 40 40 40 41 41 41 6 ...
##  $ Index.Trafic      : Factor w/ 251 levels "","0,54","0,57",..: 173 177 180 153 152 151 223 218 214 248 ...
##  $ Index.TC          : num  86.4 86.4 86.5 88.5 88.5 ...
##  $ Var.SP95..        : Factor w/ 69 levels "","-0,54%","-0,65%",..: 1 1 1 57 57 57 43 43 43 65 ...
##  $ Var.Trafic..      : Factor w/ 228 levels "","-0,01%","-0,02%",..: 1 1 1 19 36 44 227 223 217 164 ...
##  $ Var.TC..          : Factor w/ 174 levels "","-0,03%","-0,90%",..: 1 1 1 109 113 110 163 162 165 141 ...
##  $ Niveau.Prix.SP95  : Factor w/ 5 levels "","Bas","Élevé",..: 2 2 2 2 2 2 2 2 2 4 ...
##  $ Niveau.Prix.Gazole: Factor w/ 5 levels "","Bas","Élevé",..: 2 2 2 2 2 2 2 2 2 2 ...
##  $ Var.Gazole..      : Factor w/ 73 levels "","-0,56%","-0,68%",..: 1 1 1 58 58 58 43 43 43 64 ...
##  $ Index.Gazole      : num  95.7 95.7 95.7 97.8 97.8 ...
##  $ X                 : logi  NA NA NA NA NA NA ...
##  $ X.1               : Factor w/ 4 levels ""," ","GLOBALE ",..: 1 1 4 1 3 1 2 1 1 1 ...
##  $ X.2               : Factor w/ 7 levels "","1,521666667",..: 1 1 1 5 2 1 1 7 6 3 ...
##  $ X.3               : Factor w/ 6 levels "","10032,69444",..: 1 1 1 6 2 1 1 1 4 3 ...
##  $ X.4               : Factor w/ 6 levels "","10419,5","1593,583333",..: 1 1 1 6 5 1 1 1 2 3 ...
##  $ X.5               : Factor w/ 3 levels "","1,441666667",..: 1 1 1 3 2 1 1 1 1 1 ...
\end{verbatim}

L'analyse de la structure des données (str) montre que plusieurs
variables numériques ont été importées comme facteurs. C'est notamment
le cas des variables de prix, de trafic et de variations en pourcentage,
probablement en raison de la présence de symboles tels que €, \% et des
virgules décimales françaises. De plus, la variable Date a également été
importée comme facteur alors qu'un format de type date serait plus
approprié pour une analyse temporelle. Ces observations indiquent la
nécessité d'un nettoyage et d'une conversion des types avant les
traitements statistiques.

\subsubsection{nettoyage de donnés}\label{nettoyage-de-donnuxe9s}

\begin{Shaded}
\begin{Highlighting}[]
\CommentTok{\# Copie de travail}
\NormalTok{df }\OtherTok{\textless{}{-}}\NormalTok{ df\_raw}
\end{Highlighting}
\end{Shaded}

\begin{Shaded}
\begin{Highlighting}[]
\CommentTok{\# supression de colonnes de calculs}
\NormalTok{df }\OtherTok{\textless{}{-}} \FunctionTok{subset}\NormalTok{(}
\NormalTok{  df,}
  \AttributeTok{select =} \SpecialCharTok{{-}}\FunctionTok{c}\NormalTok{(X, X}\FloatTok{.1}\NormalTok{, X}\FloatTok{.2}\NormalTok{, X}\FloatTok{.3}\NormalTok{, X}\FloatTok{.4}\NormalTok{, X}\FloatTok{.5}\NormalTok{)}
\NormalTok{)}
\end{Highlighting}
\end{Shaded}

\begin{Shaded}
\begin{Highlighting}[]
\FunctionTok{str}\NormalTok{(df)}
\end{Highlighting}
\end{Shaded}

\begin{verbatim}
## 'data.frame':    274 obs. of  16 variables:
##  $ Date              : Factor w/ 93 levels "","2019-01","2019-02",..: 2 2 2 3 3 3 4 4 4 5 ...
##  $ Zone              : Factor w/ 8 levels "","0,41","0,50",..: 8 6 7 8 6 7 8 6 7 8 ...
##  $ Prix_SP95         : Factor w/ 44 levels "","1,250 €","1,300 €",..: 9 9 9 10 10 10 11 11 11 15 ...
##  $ Prix.Gazole       : Factor w/ 50 levels "","0,38","0,47",..: 14 14 14 16 16 16 18 18 18 21 ...
##  $ Trafic.Routier    : Factor w/ 260 levels "","10001","10002",..: 96 22 198 82 18 188 112 42 212 121 ...
##  $ Voyageurs.TC      : Factor w/ 231 levels "","10073","10172",..: 200 55 129 207 61 133 222 76 142 3 ...
##  $ Index.SP95        : Factor w/ 43 levels "","0,62","0,64",..: 39 39 39 40 40 40 41 41 41 6 ...
##  $ Index.Trafic      : Factor w/ 251 levels "","0,54","0,57",..: 173 177 180 153 152 151 223 218 214 248 ...
##  $ Index.TC          : num  86.4 86.4 86.5 88.5 88.5 ...
##  $ Var.SP95..        : Factor w/ 69 levels "","-0,54%","-0,65%",..: 1 1 1 57 57 57 43 43 43 65 ...
##  $ Var.Trafic..      : Factor w/ 228 levels "","-0,01%","-0,02%",..: 1 1 1 19 36 44 227 223 217 164 ...
##  $ Var.TC..          : Factor w/ 174 levels "","-0,03%","-0,90%",..: 1 1 1 109 113 110 163 162 165 141 ...
##  $ Niveau.Prix.SP95  : Factor w/ 5 levels "","Bas","Élevé",..: 2 2 2 2 2 2 2 2 2 4 ...
##  $ Niveau.Prix.Gazole: Factor w/ 5 levels "","Bas","Élevé",..: 2 2 2 2 2 2 2 2 2 2 ...
##  $ Var.Gazole..      : Factor w/ 73 levels "","-0,56%","-0,68%",..: 1 1 1 58 58 58 43 43 43 64 ...
##  $ Index.Gazole      : num  95.7 95.7 95.7 97.8 97.8 ...
\end{verbatim}

\begin{Shaded}
\begin{Highlighting}[]
\FunctionTok{summary}\NormalTok{(df)}
\end{Highlighting}
\end{Shaded}

\begin{verbatim}
##       Date              Zone      Prix_SP95    Prix.Gazole  Trafic.Routier
##         :  5   Périurbaine:87   1,900 €: 30   1,950 €: 21          : 11   
##  2019-01:  3   Rurale     :87   1,850 €: 24   1,850 €: 18   4662   :  2   
##  2019-02:  3   Urbaine    :87   1,880 €: 18   1,800 €: 15   4762   :  2   
##  2019-03:  3              : 7   1,800 €: 15   1,900 €: 15   4929   :  2   
##  2019-04:  3   0,64       : 3   1,500 €: 12   1,450 €: 12   5170   :  2   
##  2019-05:  3   0,41       : 1   1,550 €: 12   1,750 €: 12   10001  :  1   
##  (Other):254   (Other)    : 2   (Other):163   (Other):181   (Other):254   
##   Voyageurs.TC   Index.SP95   Index.Trafic    Index.TC        Var.SP95.. 
##         :  9   124,86 : 30          : 10   Min.   :  6.50          : 16  
##  252    :  5   121,58 : 24   93,07  :  3   1st Qu.: 87.37   1,06%  : 12  
##  219    :  4   123,55 : 18   96,34  :  3   Median : 95.51   -1,05% :  9  
##  232    :  3   118,29 : 15   101,05 :  2   Mean   : 91.99   -0,54% :  6  
##  234    :  3   101,86 : 12   101,91 :  2   3rd Qu.:104.73   -1,04% :  6  
##  10596  :  2   128,15 : 12   103,22 :  2   Max.   :129.39   -1,32% :  6  
##  (Other):248   (Other):163   (Other):252   NA's   :13       (Other):219  
##   Var.Trafic..    Var.TC..     Niveau.Prix.SP95  Niveau.Prix.Gazole
##         : 16          : 16             : 13               : 13     
##  -1,58% :  3   -6,12% :  6   Bas       : 66     Bas       : 90     
##  -1,86% :  3   2,35%  :  6   Élevé     :159     Élevé     :141     
##  -2,02% :  3   4,17%  :  6   Moyen     : 30     Moyen     : 15     
##  -4,11% :  3   -1,66% :  5   Très élevé:  6     Très élevé: 15     
##  2,04%  :  3   -2,00% :  5                                         
##  (Other):243   (Other):230                                         
##   Var.Gazole..  Index.Gazole   
##         : 16   Min.   : 81.16  
##  1,56%  :  9   1st Qu.:100.58  
##  -0,68% :  6   Median :122.77  
##  -1,02% :  6   Mean   :117.26  
##  -1,03% :  6   3rd Qu.:131.79  
##  -2,56% :  6   Max.   :149.13  
##  (Other):225   NA's   :13
\end{verbatim}

\begin{Shaded}
\begin{Highlighting}[]
\CommentTok{\# conversion de la date}
\NormalTok{df}\SpecialCharTok{$}\NormalTok{Date }\OtherTok{\textless{}{-}} \FunctionTok{as.Date}\NormalTok{(}
  \FunctionTok{paste0}\NormalTok{(}\FunctionTok{as.character}\NormalTok{(df}\SpecialCharTok{$}\NormalTok{Date), }\StringTok{"{-}01"}\NormalTok{),}
  \AttributeTok{format =} \StringTok{"\%Y{-}\%m{-}\%d"}
\NormalTok{)}
\end{Highlighting}
\end{Shaded}

nous avons ensuite nettoye la variable \texttt{zone}, car certaines
valeurs incorrectes avaient ete importees lors de la lecture du fichier
csv. seules les trois categories utiles a l'analyse ont ete conservees :
\texttt{urbaine}, \texttt{periurbaine} et \texttt{rurale}. cette etape
permet d'eviter des erreurs dans les analyses statistiques et les
visualisations.

\begin{Shaded}
\begin{Highlighting}[]
\CommentTok{\# nettoyage de la variable zone}
\NormalTok{df}\SpecialCharTok{$}\NormalTok{Zone }\OtherTok{\textless{}{-}} \FunctionTok{as.character}\NormalTok{(df}\SpecialCharTok{$}\NormalTok{Zone)}

\NormalTok{df }\OtherTok{\textless{}{-}} \FunctionTok{subset}\NormalTok{(}
\NormalTok{  df,}
\NormalTok{  Zone }\SpecialCharTok{\%in\%} \FunctionTok{c}\NormalTok{(}\StringTok{"Urbaine"}\NormalTok{, }\StringTok{"Périurbaine"}\NormalTok{, }\StringTok{"Rurale"}\NormalTok{)}
\NormalTok{)}

\NormalTok{df}\SpecialCharTok{$}\NormalTok{Zone }\OtherTok{\textless{}{-}} \FunctionTok{factor}\NormalTok{(}
\NormalTok{  df}\SpecialCharTok{$}\NormalTok{Zone,}
  \AttributeTok{levels =} \FunctionTok{c}\NormalTok{(}\StringTok{"Urbaine"}\NormalTok{, }\StringTok{"Périurbaine"}\NormalTok{, }\StringTok{"Rurale"}\NormalTok{)}
\NormalTok{)}
\end{Highlighting}
\end{Shaded}

\begin{Shaded}
\begin{Highlighting}[]
\CommentTok{\# function de conversion numerique}
\NormalTok{clean\_numeric }\OtherTok{\textless{}{-}} \ControlFlowTok{function}\NormalTok{(x)\{}

\NormalTok{  x }\OtherTok{\textless{}{-}} \FunctionTok{as.character}\NormalTok{(x)}

\NormalTok{  x }\OtherTok{\textless{}{-}} \FunctionTok{gsub}\NormalTok{(}\StringTok{"€"}\NormalTok{, }\StringTok{""}\NormalTok{, x)}
\NormalTok{  x }\OtherTok{\textless{}{-}} \FunctionTok{gsub}\NormalTok{(}\StringTok{"\%"}\NormalTok{, }\StringTok{""}\NormalTok{, x)}
\NormalTok{  x }\OtherTok{\textless{}{-}} \FunctionTok{gsub}\NormalTok{(}\StringTok{","}\NormalTok{, }\StringTok{"."}\NormalTok{, x)}

\NormalTok{  x }\OtherTok{\textless{}{-}} \FunctionTok{trimws}\NormalTok{(x)}

  \FunctionTok{as.numeric}\NormalTok{(x)}
\NormalTok{\}}
\end{Highlighting}
\end{Shaded}

\begin{Shaded}
\begin{Highlighting}[]
\CommentTok{\#variables numerique}
\NormalTok{vars\_numeric }\OtherTok{\textless{}{-}} \FunctionTok{c}\NormalTok{(}
  \StringTok{"Prix\_SP95"}\NormalTok{,}
  \StringTok{"Prix.Gazole"}\NormalTok{,}
  \StringTok{"Trafic.Routier"}\NormalTok{,}
  \StringTok{"Voyageurs.TC"}\NormalTok{,}
  \StringTok{"Index.SP95"}\NormalTok{,}
  \StringTok{"Index.Trafic"}\NormalTok{,}
  \StringTok{"Var.SP95.."}\NormalTok{,}
  \StringTok{"Var.Trafic.."}\NormalTok{,}
  \StringTok{"Var.TC.."}\NormalTok{,}
  \StringTok{"Var.Gazole.."}
\NormalTok{)}
\end{Highlighting}
\end{Shaded}

\begin{Shaded}
\begin{Highlighting}[]
\CommentTok{\# appliquer la fonciton numerique}
\NormalTok{df[vars\_numeric] }\OtherTok{\textless{}{-}} \FunctionTok{lapply}\NormalTok{(}
\NormalTok{  df[vars\_numeric],}
\NormalTok{  clean\_numeric}
\NormalTok{)}
\end{Highlighting}
\end{Shaded}

\begin{Shaded}
\begin{Highlighting}[]
\CommentTok{\# verification finale}
\FunctionTok{str}\NormalTok{(df)}
\end{Highlighting}
\end{Shaded}

\begin{verbatim}
## 'data.frame':    261 obs. of  16 variables:
##  $ Date              : Date, format: "2019-01-01" "2019-01-01" ...
##  $ Zone              : Factor w/ 3 levels "Urbaine","Périurbaine",..: 1 2 3 1 2 3 1 2 3 1 ...
##  $ Prix_SP95         : num  1.45 1.45 1.45 1.48 1.48 1.48 1.5 1.5 1.5 1.55 ...
##  $ Prix.Gazole       : num  1.38 1.38 1.38 1.41 1.41 1.41 1.43 1.43 1.43 1.47 ...
##  $ Trafic.Routier    : num  12562 10624 4760 12328 10409 ...
##  $ Voyageurs.TC      : num  9007 1377 212 9219 1410 ...
##  $ Index.SP95        : num  95.3 95.3 95.3 97.3 97.3 ...
##  $ Index.Trafic      : num  92.8 92.9 93 91 91 ...
##  $ Index.TC          : num  86.4 86.4 86.5 88.5 88.5 ...
##  $ Var.SP95..        : num  NA NA NA 2.07 2.07 2.07 1.35 1.35 1.35 3.33 ...
##  $ Var.Trafic..      : num  NA NA NA -1.86 -2.02 -2.18 6.96 6.74 6.53 2.34 ...
##  $ Var.TC..          : num  NA NA NA 2.35 2.4 2.36 6.89 6.88 6.91 3.23 ...
##  $ Niveau.Prix.SP95  : Factor w/ 5 levels "","Bas","Élevé",..: 2 2 2 2 2 2 2 2 2 4 ...
##  $ Niveau.Prix.Gazole: Factor w/ 5 levels "","Bas","Élevé",..: 2 2 2 2 2 2 2 2 2 2 ...
##  $ Var.Gazole..      : num  NA NA NA 2.17 2.17 2.17 1.42 1.42 1.42 2.8 ...
##  $ Index.Gazole      : num  95.7 95.7 95.7 97.8 97.8 ...
\end{verbatim}

\begin{Shaded}
\begin{Highlighting}[]
\FunctionTok{summary}\NormalTok{(df)}
\end{Highlighting}
\end{Shaded}

\begin{verbatim}
##       Date                     Zone      Prix_SP95      Prix.Gazole  
##  Min.   :2019-01-01   Urbaine    :87   Min.   :1.250   Min.   :1.17  
##  1st Qu.:2020-10-01   Périurbaine:87   1st Qu.:1.540   1st Qu.:1.45  
##  Median :2022-08-01   Rurale     :87   Median :1.820   Median :1.77  
##  Mean   :2022-08-01                    Mean   :1.728   Mean   :1.69  
##  3rd Qu.:2024-06-01                    3rd Qu.:1.890   3rd Qu.:1.90  
##  Max.   :2026-03-01                    Max.   :2.100   Max.   :2.15  
##                                                                      
##  Trafic.Routier   Voyageurs.TC     Index.SP95      Index.Trafic   
##  Min.   :  766   Min.   :   16   Min.   : 82.15   Min.   : 13.95  
##  1st Qu.: 4998   1st Qu.:  252   1st Qu.:101.20   1st Qu.: 89.76  
##  Median :10644   Median : 1510   Median :119.61   Median : 95.63  
##  Mean   : 9319   Mean   : 3758   Mean   :113.54   Mean   : 92.46  
##  3rd Qu.:12447   3rd Qu.: 9104   3rd Qu.:124.21   3rd Qu.:101.22  
##  Max.   :16008   Max.   :13467   Max.   :138.01   Max.   :118.19  
##                                                                   
##     Index.TC        Var.SP95..      Var.Trafic..        Var.TC..      
##  Min.   :  6.50   Min.   :-8.780   Min.   :-75.360   Min.   :-81.660  
##  1st Qu.: 87.37   1st Qu.:-1.270   1st Qu.: -3.958   1st Qu.: -4.370  
##  Median : 95.51   Median : 1.075   Median : -1.560   Median :  2.335  
##  Mean   : 91.99   Mean   : 0.357   Mean   :  2.449   Mean   :  4.106  
##  3rd Qu.:104.73   3rd Qu.: 1.800   3rd Qu.:  3.938   3rd Qu.:  4.390  
##  Max.   :129.39   Max.   :10.530   Max.   :208.700   Max.   :245.160  
##                   NA's   :3        NA's   :3         NA's   :3        
##    Niveau.Prix.SP95  Niveau.Prix.Gazole  Var.Gazole..       Index.Gazole   
##            :  0               :  0      Min.   :-11.6300   Min.   : 81.16  
##  Bas       : 66     Bas       : 90      1st Qu.: -1.1300   1st Qu.:100.58  
##  Élevé     :159     Élevé     :141      Median :  1.0950   Median :122.77  
##  Moyen     : 30     Moyen     : 15      Mean   :  0.5342   Mean   :117.26  
##  Très élevé:  6     Très élevé: 15      3rd Qu.:  2.1100   3rd Qu.:131.79  
##                                         Max.   : 19.4400   Max.   :149.13  
##                                         NA's   :3
\end{verbatim}

\subsubsection{créer df differents pour chaque
besoin}\label{cruxe9er-df-differents-pour-chaque-besoin}

\begin{Shaded}
\begin{Highlighting}[]
\CommentTok{\# creation du dataset pour les analyses en variations}
\NormalTok{df\_var }\OtherTok{\textless{}{-}} \FunctionTok{subset}\NormalTok{(}
\NormalTok{  df,}
  \SpecialCharTok{!}\FunctionTok{is.na}\NormalTok{(Var.SP95..) }\SpecialCharTok{\&}
  \SpecialCharTok{!}\FunctionTok{is.na}\NormalTok{(Var.Trafic..) }\SpecialCharTok{\&}
  \SpecialCharTok{!}\FunctionTok{is.na}\NormalTok{(Var.TC..)}
\NormalTok{)}
\end{Highlighting}
\end{Shaded}

\begin{Shaded}
\begin{Highlighting}[]
\CommentTok{\# creation du dataset pour les analyses en niveaux absolus}
\NormalTok{df\_abs }\OtherTok{\textless{}{-}} \FunctionTok{subset}\NormalTok{(}
\NormalTok{  df,}
  \SpecialCharTok{!}\FunctionTok{is.na}\NormalTok{(Prix\_SP95) }\SpecialCharTok{\&}
  \SpecialCharTok{!}\FunctionTok{is.na}\NormalTok{(Trafic.Routier) }\SpecialCharTok{\&}
\NormalTok{  Trafic.Routier }\SpecialCharTok{\textgreater{}} \DecValTok{0} \SpecialCharTok{\&}
\NormalTok{  Prix\_SP95 }\SpecialCharTok{\textgreater{}} \DecValTok{0}
\NormalTok{)}
\end{Highlighting}
\end{Shaded}

\begin{Shaded}
\begin{Highlighting}[]
\NormalTok{couleurs\_zone }\OtherTok{\textless{}{-}} \FunctionTok{c}\NormalTok{(}
  \StringTok{"Urbaine"} \OtherTok{=} \StringTok{"\#2196F3"}\NormalTok{,}
  \StringTok{"Périurbaine"} \OtherTok{=} \StringTok{"\#FF9800"}\NormalTok{,}
  \StringTok{"Rurale"} \OtherTok{=} \StringTok{"\#4CAF50"}
\NormalTok{)}
\end{Highlighting}
\end{Shaded}

\begin{Shaded}
\begin{Highlighting}[]
\FunctionTok{cat}\NormalTok{(}\FunctionTok{sprintf}\NormalTok{(}
  \StringTok{"donnees chargees : \%d obs. (df\_var) | \%d obs. (df\_abs)}\SpecialCharTok{\textbackslash{}n}\StringTok{"}\NormalTok{,}
  \FunctionTok{nrow}\NormalTok{(df\_var),}
  \FunctionTok{nrow}\NormalTok{(df\_abs)}
\NormalTok{))}
\end{Highlighting}
\end{Shaded}

\begin{verbatim}
## donnees chargees : 258 obs. (df_var) | 261 obs. (df_abs)
\end{verbatim}

\begin{Shaded}
\begin{Highlighting}[]
\FunctionTok{cat}\NormalTok{(}\FunctionTok{sprintf}\NormalTok{(}
  \StringTok{"periode : \%s {-}\textgreater{} \%s | zones : \%s}\SpecialCharTok{\textbackslash{}n}\StringTok{"}\NormalTok{,}
  \FunctionTok{format}\NormalTok{(}\FunctionTok{min}\NormalTok{(df\_abs}\SpecialCharTok{$}\NormalTok{Date), }\StringTok{"\%Y{-}\%m"}\NormalTok{),}
  \FunctionTok{format}\NormalTok{(}\FunctionTok{max}\NormalTok{(df\_abs}\SpecialCharTok{$}\NormalTok{Date), }\StringTok{"\%Y{-}\%m"}\NormalTok{),}
  \FunctionTok{paste}\NormalTok{(}\FunctionTok{levels}\NormalTok{(df\_abs}\SpecialCharTok{$}\NormalTok{Zone), }\AttributeTok{collapse =} \StringTok{", "}\NormalTok{)}
\NormalTok{))}
\end{Highlighting}
\end{Shaded}

\begin{verbatim}
## periode : 2019-01 -> 2026-03 | zones : Urbaine, Périurbaine, Rurale
\end{verbatim}

\subsubsection{creation de df global}\label{creation-de-df-global}

\begin{Shaded}
\begin{Highlighting}[]
\NormalTok{df\_glob }\OtherTok{\textless{}{-}}\NormalTok{ df\_abs }\SpecialCharTok{\%\textgreater{}\%}
  \FunctionTok{group\_by}\NormalTok{(Date) }\SpecialCharTok{\%\textgreater{}\%}
  \FunctionTok{summarise}\NormalTok{(}
    \AttributeTok{Prix\_SP95 =} \FunctionTok{mean}\NormalTok{(Prix\_SP95),}
    \AttributeTok{Trafic    =} \FunctionTok{mean}\NormalTok{(Trafic.Routier),}
    \AttributeTok{Voyageurs =} \FunctionTok{mean}\NormalTok{(Voyageurs.TC),}
    \AttributeTok{.groups   =} \StringTok{"drop"}
\NormalTok{  )}
\end{Highlighting}
\end{Shaded}

\begin{Shaded}
\begin{Highlighting}[]
\CommentTok{\# calcul du coefficient de mise a l\textquotesingle{}echelle}

\NormalTok{coef\_ax }\OtherTok{\textless{}{-}} \FunctionTok{max}\NormalTok{(df\_glob}\SpecialCharTok{$}\NormalTok{Trafic) }\SpecialCharTok{/} \FunctionTok{max}\NormalTok{(df\_glob}\SpecialCharTok{$}\NormalTok{Prix\_SP95)}


\CommentTok{\# creation du graphique temporel}

\NormalTok{g1 }\OtherTok{\textless{}{-}} \FunctionTok{ggplot}\NormalTok{(df\_glob, }\FunctionTok{aes}\NormalTok{(}\AttributeTok{x =}\NormalTok{ Date)) }\SpecialCharTok{+}

  \CommentTok{\# courbe du trafic routier}

  \FunctionTok{geom\_line}\NormalTok{(}
    \FunctionTok{aes}\NormalTok{(}
      \AttributeTok{y =}\NormalTok{ Trafic,}
      \AttributeTok{colour =} \StringTok{"Trafic routier"}
\NormalTok{    ),}
    \AttributeTok{linewidth =} \DecValTok{1}
\NormalTok{  ) }\SpecialCharTok{+}

  \CommentTok{\# courbe du prix SP95}

  \FunctionTok{geom\_line}\NormalTok{(}
    \FunctionTok{aes}\NormalTok{(}
      \AttributeTok{y =}\NormalTok{ Prix\_SP95 }\SpecialCharTok{*}\NormalTok{ coef\_ax,}
      \AttributeTok{colour =} \StringTok{"Prix SP95"}
\NormalTok{    ),}
    \AttributeTok{linewidth =} \DecValTok{1}\NormalTok{,}
    \AttributeTok{linetype =} \StringTok{"dashed"}
\NormalTok{  ) }\SpecialCharTok{+}

  \CommentTok{\# ajout des axes}

  \FunctionTok{scale\_y\_continuous}\NormalTok{(}
    \AttributeTok{name =} \StringTok{"Trafic routier (veh.)"}\NormalTok{,}

    \AttributeTok{sec.axis =} \FunctionTok{sec\_axis}\NormalTok{(}
      \SpecialCharTok{\textasciitilde{}}\NormalTok{ . }\SpecialCharTok{/}\NormalTok{ coef\_ax,}
      \AttributeTok{name =} \StringTok{"Prix SP95 (€/L)"}
\NormalTok{    )}
\NormalTok{  ) }\SpecialCharTok{+}

  \CommentTok{\# couleurs des courbes}

  \FunctionTok{scale\_colour\_manual}\NormalTok{(}
    \AttributeTok{values =} \FunctionTok{c}\NormalTok{(}
      \StringTok{"Trafic routier"} \OtherTok{=} \StringTok{"\#2196F3"}\NormalTok{,}
      \StringTok{"Prix SP95"} \OtherTok{=} \StringTok{"\#F44336"}
\NormalTok{    )}
\NormalTok{  ) }\SpecialCharTok{+}

  \CommentTok{\# affichage des annees}

  \FunctionTok{scale\_x\_date}\NormalTok{(}
    \AttributeTok{date\_breaks =} \StringTok{"1 year"}\NormalTok{,}
    \AttributeTok{date\_labels =} \StringTok{"\%Y"}
\NormalTok{  ) }\SpecialCharTok{+}

  \CommentTok{\# titres}

  \FunctionTok{labs}\NormalTok{(}
    \AttributeTok{title =} \StringTok{"G1 — Evolution du prix SP95 et du trafic routier (2019{-}2026)"}\NormalTok{,}
    \AttributeTok{subtitle =} \StringTok{"Donnees agregées toutes zones confondues"}
\NormalTok{  ) }\SpecialCharTok{+}

  \CommentTok{\# theme general}

  \FunctionTok{theme\_minimal}\NormalTok{(}\AttributeTok{base\_size =} \DecValTok{12}\NormalTok{) }\SpecialCharTok{+}

  \FunctionTok{theme}\NormalTok{(}
    \AttributeTok{legend.position =} \StringTok{"top"}
\NormalTok{  )}


\CommentTok{\# affichage du graphique}

\NormalTok{g1}
\end{Highlighting}
\end{Shaded}

\pandocbounded{\includegraphics[keepaspectratio]{price_elasticity_gasoline_demand_corrigé_files/figure-latex/unnamed-chunk-18-1.pdf}}

\begin{Shaded}
\begin{Highlighting}[]
\CommentTok{\# sauvegarde du graphique dans le dossier du projet}

\FunctionTok{ggsave}\NormalTok{(}
  \AttributeTok{filename =} \StringTok{"g1\_prix\_trafic.png"}\NormalTok{,}
  \AttributeTok{plot =}\NormalTok{ g1,}
  \AttributeTok{width =} \DecValTok{10}\NormalTok{,}
  \AttributeTok{height =} \DecValTok{6}\NormalTok{,}
  \AttributeTok{dpi =} \DecValTok{300}
\NormalTok{)}
\end{Highlighting}
\end{Shaded}

Le graphique présente l'évolution conjointe du trafic routier et du prix
du SP95 entre 2019 et 2026. La courbe bleue représente le trafic
routier, tandis que la courbe rouge en pointillés correspond au prix du
carburant SP95.

L'élément le plus marquant du graphique est la chute brutale du trafic
en 2020. Le nombre de véhicules diminue fortement, atteignant un niveau
exceptionnellement bas, alors que le prix du carburant ne baisse que de
manière modérée. Cette rupture correspond clairement à la crise
sanitaire de la COVID-19, marquée par les confinements, les restrictions
de circulation et la réduction générale de la mobilité. Cela suggère que
la baisse du trafic n'a pas été provoquée principalement par le prix du
carburant, mais plutôt par un choc externe lié au contexte sanitaire.

À partir de 2021, le trafic routier retrouve rapidement des niveaux
proches de ceux observés avant la pandémie, malgré une forte
augmentation du prix du SP95. Cette observation est particulièrement
intéressante dans le cadre de l'étude de l'élasticité-prix de la
demande. En théorie, une hausse des prix devrait entraîner une
diminution de la demande. Pourtant, dans ce cas, le trafic reste
relativement élevé même lorsque les prix augmentent. Cela peut indiquer
que la demande de mobilité routière est relativement inélastique à moyen
terme.

Le graphique met également en évidence une forte hausse du prix du SP95
entre 2021 et 2022, avec un niveau proche de 2 €/L. Cette évolution peut
être associée au contexte de crise énergétique européenne, à la guerre
en Ukraine et à l'augmentation générale des prix de l'énergie.

Enfin, après 2022, le trafic présente des variations régulières au cours
du temps, avec des phases de hausse et de baisse relativement cycliques.
Ce comportement peut refléter des effets de saisonnalité, notamment liés
aux périodes de vacances et aux variations saisonnières de la mobilité.

Dans l'ensemble, la relation visuelle entre les deux variables ne montre
pas une corrélation négative immédiate et très forte. Les deux courbes
n'évoluent pas systématiquement dans des directions opposées, ce qui
pourrait indiquer soit une faible élasticité-prix, soit des effets
différés dans le temps.

\subsubsection{SQ1 --- seuil de rupture
psychologique}\label{sq1-seuil-de-rupture-psychologique}

a partir de quel niveau de prix observe-t-on une baisse significative du
trafic routier ?

\begin{Shaded}
\begin{Highlighting}[]
\CommentTok{\# creation de classes de prix du SP95}
\CommentTok{\# et calcul de l\textquotesingle{}elasticite pour chaque observation}

\NormalTok{df\_seuil }\OtherTok{\textless{}{-}}\NormalTok{ df\_var }\SpecialCharTok{\%\textgreater{}\%}

  \FunctionTok{mutate}\NormalTok{(}

    \CommentTok{\# creation de 5 intervalles de prix}

    \AttributeTok{Faix\_Prix =} \FunctionTok{case\_when}\NormalTok{(}

\NormalTok{      Prix\_SP95 }\SpecialCharTok{\textless{}} \FloatTok{1.50} \SpecialCharTok{\textasciitilde{}} \StringTok{"F1: \textless{}1.50€"}\NormalTok{,}

\NormalTok{      Prix\_SP95 }\SpecialCharTok{\textgreater{}=} \FloatTok{1.50} \SpecialCharTok{\&}\NormalTok{ Prix\_SP95 }\SpecialCharTok{\textless{}} \FloatTok{1.70} \SpecialCharTok{\textasciitilde{}} \StringTok{"F2: 1.50{-}1.70€"}\NormalTok{,}

\NormalTok{      Prix\_SP95 }\SpecialCharTok{\textgreater{}=} \FloatTok{1.70} \SpecialCharTok{\&}\NormalTok{ Prix\_SP95 }\SpecialCharTok{\textless{}} \FloatTok{1.90} \SpecialCharTok{\textasciitilde{}} \StringTok{"F3: 1.70{-}1.90€"}\NormalTok{,}

\NormalTok{      Prix\_SP95 }\SpecialCharTok{\textgreater{}=} \FloatTok{1.90} \SpecialCharTok{\&}\NormalTok{ Prix\_SP95 }\SpecialCharTok{\textless{}} \FloatTok{2.10} \SpecialCharTok{\textasciitilde{}} \StringTok{"F4: 1.90{-}2.10€"}\NormalTok{,}

      \ConstantTok{TRUE} \SpecialCharTok{\textasciitilde{}} \StringTok{"F5: ≥2.10€"}
\NormalTok{    ),}

    \CommentTok{\# organisation de l\textquotesingle{}ordre des classes}

    \AttributeTok{Faix\_Prix =} \FunctionTok{factor}\NormalTok{(}
\NormalTok{      Faix\_Prix,}

      \AttributeTok{levels =} \FunctionTok{c}\NormalTok{(}
        \StringTok{"F1: \textless{}1.50€"}\NormalTok{,}
        \StringTok{"F2: 1.50{-}1.70€"}\NormalTok{,}
        \StringTok{"F3: 1.70{-}1.90€"}\NormalTok{,}
        \StringTok{"F4: 1.90{-}2.10€"}\NormalTok{,}
        \StringTok{"F5: ≥2.10€"}
\NormalTok{      )}
\NormalTok{    ),}

    \CommentTok{\# calcul de l\textquotesingle{}elasticite{-}prix}
    \CommentTok{\# epsilon = variation du trafic / variation du prix}

    \AttributeTok{epsilon\_SP95 =}\NormalTok{ Var.Trafic.. }\SpecialCharTok{/}\NormalTok{ Var.SP95..}
\NormalTok{  ) }\SpecialCharTok{\%\textgreater{}\%}

  \CommentTok{\# suppression des valeurs extremes}
  \CommentTok{\# afin d\textquotesingle{}eviter des elasticites irreales}

  \FunctionTok{filter}\NormalTok{(}\FunctionTok{abs}\NormalTok{(epsilon\_SP95) }\SpecialCharTok{\textless{}} \DecValTok{50}\NormalTok{)}


\CommentTok{\# calcul des statistiques descriptives}
\CommentTok{\# par zone et par classe de prix}

\NormalTok{seuil\_table }\OtherTok{\textless{}{-}}\NormalTok{ df\_seuil }\SpecialCharTok{\%\textgreater{}\%}

  \FunctionTok{group\_by}\NormalTok{(Zone, Faix\_Prix) }\SpecialCharTok{\%\textgreater{}\%}

  \FunctionTok{summarise}\NormalTok{(}

    \CommentTok{\# nombre d\textquotesingle{}observations}

    \AttributeTok{n\_obs =} \FunctionTok{n}\NormalTok{(),}

    \CommentTok{\# prix moyen dans chaque classe}

    \AttributeTok{Prix\_moy =} \FunctionTok{mean}\NormalTok{(Prix\_SP95),}

    \CommentTok{\# elasticite moyenne}

    \AttributeTok{eps\_moy =} \FunctionTok{mean}\NormalTok{(epsilon\_SP95, }\AttributeTok{na.rm =} \ConstantTok{TRUE}\NormalTok{),}

    \CommentTok{\# elasticite mediane}

    \AttributeTok{eps\_med =} \FunctionTok{median}\NormalTok{(epsilon\_SP95, }\AttributeTok{na.rm =} \ConstantTok{TRUE}\NormalTok{),}

    \CommentTok{\# ecart{-}type de l\textquotesingle{}elasticite}

    \AttributeTok{eps\_sd =} \FunctionTok{sd}\NormalTok{(epsilon\_SP95, }\AttributeTok{na.rm =} \ConstantTok{TRUE}\NormalTok{),}

    \CommentTok{\# variation moyenne du trafic}

    \AttributeTok{Var\_Traf =} \FunctionTok{mean}\NormalTok{(Var.Trafic.., }\AttributeTok{na.rm =} \ConstantTok{TRUE}\NormalTok{),}

    \AttributeTok{.groups =} \StringTok{"drop"}
\NormalTok{  ) }\SpecialCharTok{\%\textgreater{}\%}

  \FunctionTok{mutate}\NormalTok{(}

    \CommentTok{\# calcul des intervalles de confiance a 95 \%}

    \AttributeTok{IC\_inf =}\NormalTok{ eps\_moy }\SpecialCharTok{{-}} \FloatTok{1.96} \SpecialCharTok{*}\NormalTok{ eps\_sd }\SpecialCharTok{/} \FunctionTok{sqrt}\NormalTok{(n\_obs),}

    \AttributeTok{IC\_sup =}\NormalTok{ eps\_moy }\SpecialCharTok{+} \FloatTok{1.96} \SpecialCharTok{*}\NormalTok{ eps\_sd }\SpecialCharTok{/} \FunctionTok{sqrt}\NormalTok{(n\_obs)}
\NormalTok{  )}


\CommentTok{\# detection du seuil de rupture}
\CommentTok{\# on identifie la plus forte variation d\textquotesingle{}elasticite}
\CommentTok{\# entre deux classes de prix consecutives}

\NormalTok{seuil\_rupture }\OtherTok{\textless{}{-}}\NormalTok{ seuil\_table }\SpecialCharTok{\%\textgreater{}\%}

  \FunctionTok{group\_by}\NormalTok{(Zone) }\SpecialCharTok{\%\textgreater{}\%}

  \FunctionTok{arrange}\NormalTok{(Faix\_Prix) }\SpecialCharTok{\%\textgreater{}\%}

  \FunctionTok{mutate}\NormalTok{(}

    \CommentTok{\# difference entre deux elasticites consecutives}

    \AttributeTok{delta\_eps =}\NormalTok{ eps\_moy }\SpecialCharTok{{-}} \FunctionTok{lag}\NormalTok{(eps\_moy),}

    \CommentTok{\# identification d\textquotesingle{}une rupture importante}

    \AttributeTok{Rupture =} \FunctionTok{abs}\NormalTok{(delta\_eps) }\SpecialCharTok{\textgreater{}} \FloatTok{0.15}
\NormalTok{  ) }\SpecialCharTok{\%\textgreater{}\%}

  \FunctionTok{filter}\NormalTok{(}\SpecialCharTok{!}\FunctionTok{is.na}\NormalTok{(delta\_eps)) }\SpecialCharTok{\%\textgreater{}\%}

  \FunctionTok{summarise}\NormalTok{(}

    \CommentTok{\# classe de prix correspondant}
    \CommentTok{\# au plus grand saut d\textquotesingle{}elasticite}

    \AttributeTok{Faix\_Rupture =}\NormalTok{ Faix\_Prix[}\FunctionTok{which.max}\NormalTok{(}\FunctionTok{abs}\NormalTok{(delta\_eps))],}

    \AttributeTok{Delta\_max =} \FunctionTok{round}\NormalTok{(}\FunctionTok{max}\NormalTok{(}\FunctionTok{abs}\NormalTok{(delta\_eps), }\AttributeTok{na.rm =} \ConstantTok{TRUE}\NormalTok{), }\DecValTok{4}\NormalTok{),}

    \AttributeTok{.groups =} \StringTok{"drop"}
\NormalTok{  )}


\CommentTok{\# affichage des resultats numeriques}

\FunctionTok{cat}\NormalTok{(}\StringTok{"}\SpecialCharTok{\textbackslash{}n}\StringTok{elasticites par classe de prix et par zone :}\SpecialCharTok{\textbackslash{}n}\StringTok{"}\NormalTok{)}
\end{Highlighting}
\end{Shaded}

\begin{verbatim}
## 
## elasticites par classe de prix et par zone :
\end{verbatim}

\begin{Shaded}
\begin{Highlighting}[]
\NormalTok{seuil\_table }\SpecialCharTok{\%\textgreater{}\%}

  \FunctionTok{select}\NormalTok{(}
\NormalTok{    Zone,}
\NormalTok{    Faix\_Prix,}
\NormalTok{    n\_obs,}
\NormalTok{    Prix\_moy,}
\NormalTok{    eps\_moy,}
\NormalTok{    eps\_med,}
\NormalTok{    IC\_inf,}
\NormalTok{    IC\_sup,}
\NormalTok{    Var\_Traf}
\NormalTok{  ) }\SpecialCharTok{\%\textgreater{}\%}

  \FunctionTok{mutate}\NormalTok{(}
    \FunctionTok{across}\NormalTok{(}\FunctionTok{where}\NormalTok{(is.numeric), }\SpecialCharTok{\textasciitilde{}} \FunctionTok{round}\NormalTok{(.x, }\DecValTok{4}\NormalTok{))}
\NormalTok{  ) }\SpecialCharTok{\%\textgreater{}\%}

  \FunctionTok{print}\NormalTok{(}\AttributeTok{n =} \ConstantTok{Inf}\NormalTok{)}
\end{Highlighting}
\end{Shaded}

\begin{verbatim}
## # A tibble: 15 x 9
##    Zone        Faix_Prix   n_obs Prix_moy eps_moy eps_med IC_inf IC_sup Var_Traf
##    <fct>       <fct>       <dbl>    <dbl>   <dbl>   <dbl>  <dbl>  <dbl>    <dbl>
##  1 Urbaine     F1: <1.50€     12     1.40   5.64    2.99  -0.511 11.8     -4.07 
##  2 Urbaine     F2: 1.50-1~    18     1.56   2.99    1.31   0.424  5.56     4.08 
##  3 Urbaine     F3: 1.70-1~    33     1.83  -0.480  -0.480 -1.68   0.714   -1.07 
##  4 Urbaine     F4: 1.90-2~    20     1.92  -0.707  -0.673 -2.29   0.880    0.216
##  5 Urbaine     F5: ≥2.10€      1     2.1    1.38    1.38  NA     NA       14.6  
##  6 Périurbaine F1: <1.50€     12     1.40   5.55    2.92  -0.534 11.6     -4.10 
##  7 Périurbaine F2: 1.50-1~    18     1.56   2.88    1.19   0.356  5.41     3.93 
##  8 Périurbaine F3: 1.70-1~    33     1.83  -0.479  -0.406 -1.67   0.708   -1.01 
##  9 Périurbaine F4: 1.90-2~    20     1.92  -0.725  -0.758 -2.20   0.752    0.251
## 10 Périurbaine F5: ≥2.10€      1     2.1    0.959   0.959 NA     NA       10.1  
## 11 Rurale      F1: <1.50€     12     1.40   5.46    2.85  -0.559 11.5     -4.13 
## 12 Rurale      F2: 1.50-1~    18     1.56   2.77    1.06   0.290  5.26     3.77 
## 13 Rurale      F3: 1.70-1~    33     1.83  -0.478  -0.025 -1.66   0.705   -0.942
## 14 Rurale      F4: 1.90-2~    20     1.92  -0.743  -0.937 -2.12   0.632    0.291
## 15 Rurale      F5: ≥2.10€      1     2.1    0.511   0.511 NA     NA        5.38
\end{verbatim}

\begin{Shaded}
\begin{Highlighting}[]
\CommentTok{\# affichage du seuil detecte}

\FunctionTok{cat}\NormalTok{(}\StringTok{"}\SpecialCharTok{\textbackslash{}n}\StringTok{seuil de rupture detecte numeriquement :}\SpecialCharTok{\textbackslash{}n}\StringTok{"}\NormalTok{)}
\end{Highlighting}
\end{Shaded}

\begin{verbatim}
## 
## seuil de rupture detecte numeriquement :
\end{verbatim}

\begin{Shaded}
\begin{Highlighting}[]
\FunctionTok{print}\NormalTok{(seuil\_rupture)}
\end{Highlighting}
\end{Shaded}

\begin{verbatim}
## # A tibble: 3 x 3
##   Zone        Faix_Rupture   Delta_max
##   <fct>       <fct>              <dbl>
## 1 Urbaine     F3: 1.70-1.90€      3.47
## 2 Périurbaine F3: 1.70-1.90€      3.36
## 3 Rurale      F3: 1.70-1.90€      3.25
\end{verbatim}

\begin{Shaded}
\begin{Highlighting}[]
\CommentTok{\# creation du graphique des elasticites}

\NormalTok{g3 }\OtherTok{\textless{}{-}} \FunctionTok{ggplot}\NormalTok{(}
\NormalTok{  seuil\_table,}
  \FunctionTok{aes}\NormalTok{(}
    \AttributeTok{x =}\NormalTok{ Faix\_Prix,}
    \AttributeTok{y =}\NormalTok{ eps\_moy,}
    \AttributeTok{fill =}\NormalTok{ Zone}
\NormalTok{  )}
\NormalTok{) }\SpecialCharTok{+}

  \CommentTok{\# creation des barres}

  \FunctionTok{geom\_col}\NormalTok{(}
    \AttributeTok{position =} \StringTok{"dodge"}\NormalTok{,}
    \AttributeTok{alpha =} \FloatTok{0.85}
\NormalTok{  ) }\SpecialCharTok{+}

  \CommentTok{\# ajout des intervalles de confiance}

  \FunctionTok{geom\_errorbar}\NormalTok{(}
    \FunctionTok{aes}\NormalTok{(}
      \AttributeTok{ymin =}\NormalTok{ IC\_inf,}
      \AttributeTok{ymax =}\NormalTok{ IC\_sup}
\NormalTok{    ),}

    \AttributeTok{position =} \FunctionTok{position\_dodge}\NormalTok{(}\FloatTok{0.9}\NormalTok{),}

    \AttributeTok{width =} \FloatTok{0.25}\NormalTok{,}

    \AttributeTok{linewidth =} \FloatTok{0.7}
\NormalTok{  ) }\SpecialCharTok{+}

  \CommentTok{\# ligne horizontale correspondant a zero}

  \FunctionTok{geom\_hline}\NormalTok{(}
    \AttributeTok{yintercept =} \DecValTok{0}\NormalTok{,}
    \AttributeTok{colour =} \StringTok{"black"}
\NormalTok{  ) }\SpecialCharTok{+}

  \CommentTok{\# ligne verticale indiquant le seuil detecte}

  \FunctionTok{geom\_vline}\NormalTok{(}
    \AttributeTok{xintercept =} \FloatTok{3.5}\NormalTok{,}
    \AttributeTok{linetype =} \StringTok{"dashed"}\NormalTok{,}
    \AttributeTok{colour =} \StringTok{"\#F44336"}\NormalTok{,}
    \AttributeTok{linewidth =} \DecValTok{1}
\NormalTok{  ) }\SpecialCharTok{+}

  \CommentTok{\# annotation du seuil}

  \FunctionTok{annotate}\NormalTok{(}
    \StringTok{"text"}\NormalTok{,}

    \AttributeTok{x =} \FloatTok{3.55}\NormalTok{,}

    \AttributeTok{y =} \FunctionTok{max}\NormalTok{(seuil\_table}\SpecialCharTok{$}\NormalTok{eps\_moy, }\AttributeTok{na.rm =} \ConstantTok{TRUE}\NormalTok{) }\SpecialCharTok{*} \FloatTok{0.85}\NormalTok{,}

    \AttributeTok{label =} \StringTok{"\textless{}{-} seuil detecte"}\NormalTok{,}

    \AttributeTok{hjust =} \DecValTok{0}\NormalTok{,}

    \AttributeTok{size =} \FloatTok{3.5}\NormalTok{,}

    \AttributeTok{colour =} \StringTok{"\#F44336"}
\NormalTok{  ) }\SpecialCharTok{+}

  \CommentTok{\# couleurs selon la zone}

  \FunctionTok{scale\_fill\_manual}\NormalTok{(}
    \AttributeTok{values =}\NormalTok{ couleurs\_zone}
\NormalTok{  ) }\SpecialCharTok{+}

  \CommentTok{\# titres du graphique}

  \FunctionTok{labs}\NormalTok{(}
    \AttributeTok{title =} \StringTok{"G3 {-} SQ1 : elasticite selon les classes de prix"}\NormalTok{,}

    \AttributeTok{subtitle =} \StringTok{"Detection du seuil psychologique et intervalles de confiance a 95 \%"}\NormalTok{,}

    \AttributeTok{x =} \StringTok{"Classe de prix SP95"}\NormalTok{,}

    \AttributeTok{y =} \StringTok{"Elasticite moyenne (epsilon)"}\NormalTok{,}

    \AttributeTok{fill =} \StringTok{"Zone"}
\NormalTok{  ) }\SpecialCharTok{+}

  \CommentTok{\# mise en forme generale}

  \FunctionTok{theme\_minimal}\NormalTok{(}\AttributeTok{base\_size =} \DecValTok{12}\NormalTok{) }\SpecialCharTok{+}

  \FunctionTok{theme}\NormalTok{(}
    \AttributeTok{axis.text.x =} \FunctionTok{element\_text}\NormalTok{(}
      \AttributeTok{angle =} \DecValTok{15}\NormalTok{,}
      \AttributeTok{hjust =} \DecValTok{1}
\NormalTok{    )}
\NormalTok{  )}


\CommentTok{\# affichage du graphique}

\NormalTok{g3}
\end{Highlighting}
\end{Shaded}

\pandocbounded{\includegraphics[keepaspectratio]{price_elasticity_gasoline_demand_corrigé_files/figure-latex/unnamed-chunk-19-1.pdf}}

\begin{Shaded}
\begin{Highlighting}[]
\CommentTok{\# sauvegarde du graphique dans le dossier du projet}

\FunctionTok{ggsave}\NormalTok{(}
  \AttributeTok{filename =} \StringTok{"g3\_seuil\_rupture.png"}\NormalTok{,}
  \AttributeTok{plot =}\NormalTok{ g3,}
  \AttributeTok{width =} \DecValTok{10}\NormalTok{,}
  \AttributeTok{height =} \DecValTok{6}\NormalTok{,}
  \AttributeTok{dpi =} \DecValTok{300}
\NormalTok{)}
\end{Highlighting}
\end{Shaded}

Le point central est que le seuil de rupture apparait de maniere assez
coherente dans la classe F3 : 1.70--1.90€ pour les trois zones. Cela
signifie que, selon le critere numerique utilise, le changement le plus
important d'elasticite se produit lorsque le prix entre dans cette
tranche. En termes simples, le graphique suggere que le comportement du
trafic change plus clairement a partir d'environ 1,70€.

Cependant, plusieurs precautions d'interpretation sont necessaires.

Les valeurs d'elasticite sont tres instables dans les classes de prix
les plus basses, notamment en F1 et F2. Les intervalles de confiance
sont tres larges, ce qui montre que les moyennes observees dans ces
classes ne sont pas totalement fiables comme valeurs economiques
absolues.

Par ailleurs, la classe F5 ne contient qu'une seule observation pour
chaque zone. Les resultats associes a cette classe doivent donc etre
interpretes avec prudence et ne peuvent pas etre consideres comme
statistiquement robustes.

Le fait que les elasticites soient positives en F1 et F2 peut egalement
indiquer la presence de bruit statistique, d'effets externes ou encore
d'un contexte macroeconomique commun influencant simultanement les prix
et le trafic.

A partir de la classe F3, les elasticites deviennent proches de zero ou
negatives, ce qui correspond davantage a l'idee d'un point de rupture
psychologique.

\subsubsection{SQ2 - ELASTICITE CROISEE \& REPORT
MODAL}\label{sq2---elasticite-croisee-report-modal}

la hausse du prix corrèle-t-elle avec une augmentation des tc ?

\begin{Shaded}
\begin{Highlighting}[]
\FunctionTok{cat}\NormalTok{(}\StringTok{"}\SpecialCharTok{\textbackslash{}n}\StringTok{━━━━━━━━━━━━━━━━━━━━━━━━━━━━━━━━━━━━━━━━━━━━━━━━━━━━━━━━━}\SpecialCharTok{\textbackslash{}n}\StringTok{"}\NormalTok{)}
\end{Highlighting}
\end{Shaded}

\begin{verbatim}
## 
## ━━━━━━━━━━━━━━━━━━━━━━━━━━━━━━━━━━━━━━━━━━━━━━━━━━━━━━━━━
\end{verbatim}

\begin{Shaded}
\begin{Highlighting}[]
\FunctionTok{cat}\NormalTok{(}\StringTok{"  SQ2 {-} elasticite croisee et report modal}\SpecialCharTok{\textbackslash{}n}\StringTok{"}\NormalTok{)}
\end{Highlighting}
\end{Shaded}

\begin{verbatim}
##   SQ2 - elasticite croisee et report modal
\end{verbatim}

\begin{Shaded}
\begin{Highlighting}[]
\FunctionTok{cat}\NormalTok{(}\StringTok{"  epsilon croisee = variation TC / variation prix SP95}\SpecialCharTok{\textbackslash{}n}\StringTok{"}\NormalTok{)}
\end{Highlighting}
\end{Shaded}

\begin{verbatim}
##   epsilon croisee = variation TC / variation prix SP95
\end{verbatim}

\begin{Shaded}
\begin{Highlighting}[]
\FunctionTok{cat}\NormalTok{(}\StringTok{"━━━━━━━━━━━━━━━━━━━━━━━━━━━━━━━━━━━━━━━━━━━━━━━━━━━━━━━━━}\SpecialCharTok{\textbackslash{}n}\StringTok{"}\NormalTok{)}
\end{Highlighting}
\end{Shaded}

\begin{verbatim}
## ━━━━━━━━━━━━━━━━━━━━━━━━━━━━━━━━━━━━━━━━━━━━━━━━━━━━━━━━━
\end{verbatim}

\begin{Shaded}
\begin{Highlighting}[]
\CommentTok{\# creation du dataset pour l\textquotesingle{}analyse croisee}

\NormalTok{df\_cross }\OtherTok{\textless{}{-}}\NormalTok{ df\_var }\SpecialCharTok{\%\textgreater{}\%}

  \FunctionTok{mutate}\NormalTok{(}

    \CommentTok{\# calcul de l\textquotesingle{}elasticite croisee}
    \CommentTok{\# variation des transports en commun}
    \CommentTok{\# divisee par la variation du prix SP95}

    \AttributeTok{epsilon\_croisee =}\NormalTok{ Var.TC.. }\SpecialCharTok{/}\NormalTok{ Var.SP95..,}

    \CommentTok{\# creation des classes de prix}

    \AttributeTok{Faix\_Prix =} \FunctionTok{case\_when}\NormalTok{(}

\NormalTok{      Prix\_SP95 }\SpecialCharTok{\textless{}} \FloatTok{1.50} \SpecialCharTok{\textasciitilde{}} \StringTok{"F1: \textless{}1.50EUR"}\NormalTok{,}

\NormalTok{      Prix\_SP95 }\SpecialCharTok{\textgreater{}=} \FloatTok{1.50} \SpecialCharTok{\&}\NormalTok{ Prix\_SP95 }\SpecialCharTok{\textless{}} \FloatTok{1.70} \SpecialCharTok{\textasciitilde{}} \StringTok{"F2: 1.50{-}1.70EUR"}\NormalTok{,}

\NormalTok{      Prix\_SP95 }\SpecialCharTok{\textgreater{}=} \FloatTok{1.70} \SpecialCharTok{\&}\NormalTok{ Prix\_SP95 }\SpecialCharTok{\textless{}} \FloatTok{1.90} \SpecialCharTok{\textasciitilde{}} \StringTok{"F3: 1.70{-}1.90EUR"}\NormalTok{,}

\NormalTok{      Prix\_SP95 }\SpecialCharTok{\textgreater{}=} \FloatTok{1.90} \SpecialCharTok{\&}\NormalTok{ Prix\_SP95 }\SpecialCharTok{\textless{}} \FloatTok{2.10} \SpecialCharTok{\textasciitilde{}} \StringTok{"F4: 1.90{-}2.10EUR"}\NormalTok{,}

      \ConstantTok{TRUE} \SpecialCharTok{\textasciitilde{}} \StringTok{"F5: \textgreater{}=2.10EUR"}
\NormalTok{    ),}

    \CommentTok{\# organisation des classes dans l\textquotesingle{}ordre croissant}

    \AttributeTok{Faix\_Prix =} \FunctionTok{factor}\NormalTok{(}

\NormalTok{      Faix\_Prix,}

      \AttributeTok{levels =} \FunctionTok{c}\NormalTok{(}
        \StringTok{"F1: \textless{}1.50EUR"}\NormalTok{,}
        \StringTok{"F2: 1.50{-}1.70EUR"}\NormalTok{,}
        \StringTok{"F3: 1.70{-}1.90EUR"}\NormalTok{,}
        \StringTok{"F4: 1.90{-}2.10EUR"}\NormalTok{,}
        \StringTok{"F5: \textgreater{}=2.10EUR"}
\NormalTok{      )}
\NormalTok{    )}
\NormalTok{  ) }\SpecialCharTok{\%\textgreater{}\%}

  \CommentTok{\# suppression des valeurs extremes}

  \FunctionTok{filter}\NormalTok{(}\FunctionTok{abs}\NormalTok{(epsilon\_croisee) }\SpecialCharTok{\textless{}} \DecValTok{50}\NormalTok{)}


\CommentTok{\# calcul des statistiques par zone}

\NormalTok{cross\_zone }\OtherTok{\textless{}{-}}\NormalTok{ df\_cross }\SpecialCharTok{\%\textgreater{}\%}

  \FunctionTok{group\_by}\NormalTok{(Zone) }\SpecialCharTok{\%\textgreater{}\%}

  \FunctionTok{summarise}\NormalTok{(}

    \CommentTok{\# nombre d\textquotesingle{}observations}

    \AttributeTok{n =} \FunctionTok{n}\NormalTok{(),}

    \CommentTok{\# elasticite croisee moyenne}

    \AttributeTok{eps\_crois\_moy =} \FunctionTok{mean}\NormalTok{(epsilon\_croisee, }\AttributeTok{na.rm =} \ConstantTok{TRUE}\NormalTok{),}

    \CommentTok{\# elasticite mediane}

    \AttributeTok{eps\_crois\_med =} \FunctionTok{median}\NormalTok{(epsilon\_croisee, }\AttributeTok{na.rm =} \ConstantTok{TRUE}\NormalTok{),}

    \CommentTok{\# ecart{-}type}

    \AttributeTok{eps\_crois\_sd =} \FunctionTok{sd}\NormalTok{(epsilon\_croisee, }\AttributeTok{na.rm =} \ConstantTok{TRUE}\NormalTok{),}

    \AttributeTok{.groups =} \StringTok{"drop"}
\NormalTok{  ) }\SpecialCharTok{\%\textgreater{}\%}

  \FunctionTok{mutate}\NormalTok{(}

    \CommentTok{\# calcul des intervalles de confiance a 95 \%}

    \AttributeTok{IC\_inf =}\NormalTok{ eps\_crois\_moy }\SpecialCharTok{{-}} \FloatTok{1.96} \SpecialCharTok{*}\NormalTok{ eps\_crois\_sd }\SpecialCharTok{/} \FunctionTok{sqrt}\NormalTok{(n),}

    \AttributeTok{IC\_sup =}\NormalTok{ eps\_crois\_moy }\SpecialCharTok{+} \FloatTok{1.96} \SpecialCharTok{*}\NormalTok{ eps\_crois\_sd }\SpecialCharTok{/} \FunctionTok{sqrt}\NormalTok{(n),}

    \CommentTok{\# interpretation economique automatique}

    \AttributeTok{Interpretation =} \FunctionTok{case\_when}\NormalTok{(}

\NormalTok{      eps\_crois\_moy }\SpecialCharTok{\textgreater{}} \FloatTok{0.10} \SpecialCharTok{\textasciitilde{}}
        \StringTok{"Biens substituts {-}\textgreater{} report modal avere"}\NormalTok{,}

\NormalTok{      eps\_crois\_moy }\SpecialCharTok{\textgreater{}} \DecValTok{0} \SpecialCharTok{\textasciitilde{}}
        \StringTok{"Legere substituabilite"}\NormalTok{,}

      \ConstantTok{TRUE} \SpecialCharTok{\textasciitilde{}}
        \StringTok{"Biens independants {-}\textgreater{} pas de report modal"}
\NormalTok{    )}
\NormalTok{  )}


\CommentTok{\# affichage du tableau de synthese}

\FunctionTok{cat}\NormalTok{(}\StringTok{"}\SpecialCharTok{\textbackslash{}n}\StringTok{ tableau synthese SQ2 :}\SpecialCharTok{\textbackslash{}n}\StringTok{"}\NormalTok{)}
\end{Highlighting}
\end{Shaded}

\begin{verbatim}
## 
##  tableau synthese SQ2 :
\end{verbatim}

\begin{Shaded}
\begin{Highlighting}[]
\NormalTok{cross\_zone }\SpecialCharTok{\%\textgreater{}\%}

  \FunctionTok{select}\NormalTok{(}
\NormalTok{    Zone,}
\NormalTok{    eps\_crois\_moy,}
\NormalTok{    eps\_crois\_med,}
\NormalTok{    IC\_inf,}
\NormalTok{    IC\_sup,}
\NormalTok{    Interpretation}
\NormalTok{  ) }\SpecialCharTok{\%\textgreater{}\%}

  \FunctionTok{rename}\NormalTok{(}

    \StringTok{"epsilon croisee (moy)"} \OtherTok{=}\NormalTok{ eps\_crois\_moy,}

    \StringTok{"epsilon croisee (med)"} \OtherTok{=}\NormalTok{ eps\_crois\_med,}

    \StringTok{"IC 95\% inf"} \OtherTok{=}\NormalTok{ IC\_inf,}

    \StringTok{"IC 95\% sup"} \OtherTok{=}\NormalTok{ IC\_sup}
\NormalTok{  ) }\SpecialCharTok{\%\textgreater{}\%}

  \FunctionTok{mutate}\NormalTok{(}
    \FunctionTok{across}\NormalTok{(}\FunctionTok{where}\NormalTok{(is.numeric), }\SpecialCharTok{\textasciitilde{}} \FunctionTok{round}\NormalTok{(.x, }\DecValTok{4}\NormalTok{))}
\NormalTok{  ) }\SpecialCharTok{\%\textgreater{}\%}

  \FunctionTok{print}\NormalTok{()}
\end{Highlighting}
\end{Shaded}

\begin{verbatim}
## # A tibble: 3 x 6
##   Zone   epsilon croisee (moy~1 epsilon croisee (med~2 `IC 95% inf` `IC 95% sup`
##   <fct>                   <dbl>                  <dbl>        <dbl>        <dbl>
## 1 Urbai~                   1.36                  0.740       -0.516         3.23
## 2 Périu~                   1.35                  0.735       -0.518         3.22
## 3 Rurale                   1.33                  0.846       -0.543         3.20
## # i abbreviated names: 1: `epsilon croisee (moy)`, 2: `epsilon croisee (med)`
## # i 1 more variable: Interpretation <chr>
\end{verbatim}

\begin{Shaded}
\begin{Highlighting}[]
\CommentTok{\# analyse de l\textquotesingle{}elasticite croisee}
\CommentTok{\# selon les classes de prix}

\NormalTok{cross\_faix }\OtherTok{\textless{}{-}}\NormalTok{ df\_cross }\SpecialCharTok{\%\textgreater{}\%}

  \FunctionTok{group\_by}\NormalTok{(Zone, Faix\_Prix) }\SpecialCharTok{\%\textgreater{}\%}

  \FunctionTok{summarise}\NormalTok{(}

    \AttributeTok{eps\_croisee =} \FunctionTok{mean}\NormalTok{(}
\NormalTok{      epsilon\_croisee,}
      \AttributeTok{na.rm =} \ConstantTok{TRUE}
\NormalTok{    ),}

    \AttributeTok{.groups =} \StringTok{"drop"}
\NormalTok{  )}


\CommentTok{\# affichage des resultats par classe de prix}

\FunctionTok{cat}\NormalTok{(}\StringTok{"}\SpecialCharTok{\textbackslash{}n}\StringTok{ elasticite croisee par classe de prix :}\SpecialCharTok{\textbackslash{}n}\StringTok{"}\NormalTok{)}
\end{Highlighting}
\end{Shaded}

\begin{verbatim}
## 
##  elasticite croisee par classe de prix :
\end{verbatim}

\begin{Shaded}
\begin{Highlighting}[]
\NormalTok{cross\_faix }\SpecialCharTok{\%\textgreater{}\%}

  \FunctionTok{mutate}\NormalTok{(}
    \FunctionTok{across}\NormalTok{(}\FunctionTok{where}\NormalTok{(is.numeric), }\SpecialCharTok{\textasciitilde{}} \FunctionTok{round}\NormalTok{(.x, }\DecValTok{4}\NormalTok{))}
\NormalTok{  ) }\SpecialCharTok{\%\textgreater{}\%}

  \FunctionTok{print}\NormalTok{(}\AttributeTok{n =} \ConstantTok{Inf}\NormalTok{)}
\end{Highlighting}
\end{Shaded}

\begin{verbatim}
## # A tibble: 15 x 3
##    Zone        Faix_Prix        eps_croisee
##    <fct>       <fct>                  <dbl>
##  1 Urbaine     F1: <1.50EUR           8.40 
##  2 Urbaine     F2: 1.50-1.70EUR       4.14 
##  3 Urbaine     F3: 1.70-1.90EUR      -0.726
##  4 Urbaine     F4: 1.90-2.10EUR      -1.90 
##  5 Urbaine     F5: >=2.10EUR          0.655
##  6 Périurbaine F1: <1.50EUR           8.36 
##  7 Périurbaine F2: 1.50-1.70EUR       4.14 
##  8 Périurbaine F3: 1.70-1.90EUR      -0.723
##  9 Périurbaine F4: 1.90-2.10EUR      -1.91 
## 10 Périurbaine F5: >=2.10EUR          0.66 
## 11 Rurale      F1: <1.50EUR           8.35 
## 12 Rurale      F2: 1.50-1.70EUR       4.14 
## 13 Rurale      F3: 1.70-1.90EUR      -0.754
## 14 Rurale      F4: 1.90-2.10EUR      -1.94 
## 15 Rurale      F5: >=2.10EUR          0.650
\end{verbatim}

\begin{Shaded}
\begin{Highlighting}[]
\CommentTok{\# graphique des elasticites croisees moyennes}

\NormalTok{g4 }\OtherTok{\textless{}{-}} \FunctionTok{ggplot}\NormalTok{(}

\NormalTok{  cross\_zone,}

  \FunctionTok{aes}\NormalTok{(}
    \AttributeTok{x =}\NormalTok{ Zone,}
    \AttributeTok{y =}\NormalTok{ eps\_crois\_moy,}
    \AttributeTok{fill =}\NormalTok{ Zone}
\NormalTok{  )}
\NormalTok{) }\SpecialCharTok{+}

  \CommentTok{\# creation des barres}

  \FunctionTok{geom\_col}\NormalTok{(}
    \AttributeTok{alpha =} \FloatTok{0.85}\NormalTok{,}
    \AttributeTok{width =} \FloatTok{0.6}
\NormalTok{  ) }\SpecialCharTok{+}

  \CommentTok{\# ajout des intervalles de confiance}

  \FunctionTok{geom\_errorbar}\NormalTok{(}

    \FunctionTok{aes}\NormalTok{(}
      \AttributeTok{ymin =}\NormalTok{ IC\_inf,}
      \AttributeTok{ymax =}\NormalTok{ IC\_sup}
\NormalTok{    ),}

    \AttributeTok{width =} \FloatTok{0.2}\NormalTok{,}

    \AttributeTok{linewidth =} \FloatTok{0.8}
\NormalTok{  ) }\SpecialCharTok{+}

  \CommentTok{\# ligne de reference a zero}

  \FunctionTok{geom\_hline}\NormalTok{(}
    \AttributeTok{yintercept =} \DecValTok{0}\NormalTok{,}
    \AttributeTok{linetype =} \StringTok{"solid"}
\NormalTok{  ) }\SpecialCharTok{+}

  \CommentTok{\# couleurs des zones}

  \FunctionTok{scale\_fill\_manual}\NormalTok{(}
    \AttributeTok{values =}\NormalTok{ couleurs\_zone}
\NormalTok{  ) }\SpecialCharTok{+}

  \CommentTok{\# titres du graphique}

  \FunctionTok{labs}\NormalTok{(}

    \AttributeTok{title =} \StringTok{"G4 {-} SQ2 : elasticite croisee et report modal"}\NormalTok{,}

    \AttributeTok{subtitle =} \StringTok{"epsilon \textgreater{} 0 : substitution | epsilon proche de 0 : independance"}\NormalTok{,}

    \AttributeTok{x =} \ConstantTok{NULL}\NormalTok{,}

    \AttributeTok{y =} \StringTok{"Elasticite croisee"}\NormalTok{,}

    \AttributeTok{fill =} \ConstantTok{NULL}
\NormalTok{  ) }\SpecialCharTok{+}

  \CommentTok{\# mise en forme generale}

  \FunctionTok{theme\_minimal}\NormalTok{(}\AttributeTok{base\_size =} \DecValTok{12}\NormalTok{) }\SpecialCharTok{+}

  \FunctionTok{theme}\NormalTok{(}
    \AttributeTok{legend.position =} \StringTok{"none"}
\NormalTok{  )}


\CommentTok{\# graphique de l\textquotesingle{}evolution de l\textquotesingle{}elasticite}
\CommentTok{\# selon les classes de prix}

\NormalTok{g6 }\OtherTok{\textless{}{-}} \FunctionTok{ggplot}\NormalTok{(}

\NormalTok{  cross\_faix,}

  \FunctionTok{aes}\NormalTok{(}
    \AttributeTok{x =}\NormalTok{ Faix\_Prix,}
    \AttributeTok{y =}\NormalTok{ eps\_croisee,}
    \AttributeTok{colour =}\NormalTok{ Zone,}
    \AttributeTok{group =}\NormalTok{ Zone}
\NormalTok{  )}
\NormalTok{) }\SpecialCharTok{+}

  \CommentTok{\# courbes par zone}

  \FunctionTok{geom\_line}\NormalTok{(}
    \AttributeTok{linewidth =} \FloatTok{1.2}
\NormalTok{  ) }\SpecialCharTok{+}

  \CommentTok{\# ajout des points}

  \FunctionTok{geom\_point}\NormalTok{(}
    \AttributeTok{size =} \DecValTok{3}
\NormalTok{  ) }\SpecialCharTok{+}

  \CommentTok{\# ligne horizontale a zero}

  \FunctionTok{geom\_hline}\NormalTok{(}
    \AttributeTok{yintercept =} \DecValTok{0}\NormalTok{,}
    \AttributeTok{linetype =} \StringTok{"dashed"}
\NormalTok{  ) }\SpecialCharTok{+}

  \CommentTok{\# couleurs des zones}

  \FunctionTok{scale\_colour\_manual}\NormalTok{(}
    \AttributeTok{values =}\NormalTok{ couleurs\_zone}
\NormalTok{  ) }\SpecialCharTok{+}

  \CommentTok{\# titres du graphique}

  \FunctionTok{labs}\NormalTok{(}

    \AttributeTok{title =} \StringTok{"G6 {-} SQ2 : intensite du report modal selon le prix"}\NormalTok{,}

    \AttributeTok{subtitle =} \StringTok{"le report modal augmente{-}t{-}il lorsque le prix franchit le seuil ?"}\NormalTok{,}

    \AttributeTok{x =} \StringTok{"Classe de prix SP95"}\NormalTok{,}

    \AttributeTok{y =} \StringTok{"Elasticite croisee"}\NormalTok{,}

    \AttributeTok{colour =} \StringTok{"Zone"}
\NormalTok{  ) }\SpecialCharTok{+}

  \CommentTok{\# mise en forme generale}

  \FunctionTok{theme\_minimal}\NormalTok{(}\AttributeTok{base\_size =} \DecValTok{12}\NormalTok{) }\SpecialCharTok{+}

  \FunctionTok{theme}\NormalTok{(}
    \AttributeTok{axis.text.x =} \FunctionTok{element\_text}\NormalTok{(}
      \AttributeTok{angle =} \DecValTok{15}\NormalTok{,}
      \AttributeTok{hjust =} \DecValTok{1}
\NormalTok{    )}
\NormalTok{  )}


\CommentTok{\# affichage des graphiques}

\NormalTok{g4}
\end{Highlighting}
\end{Shaded}

\pandocbounded{\includegraphics[keepaspectratio]{price_elasticity_gasoline_demand_corrigé_files/figure-latex/unnamed-chunk-20-1.pdf}}

\begin{Shaded}
\begin{Highlighting}[]
\NormalTok{g6}
\end{Highlighting}
\end{Shaded}

\pandocbounded{\includegraphics[keepaspectratio]{price_elasticity_gasoline_demand_corrigé_files/figure-latex/unnamed-chunk-20-2.pdf}}

\begin{Shaded}
\begin{Highlighting}[]
\CommentTok{\# sauvegarde du graphique G4}

\FunctionTok{ggsave}\NormalTok{(}
  \AttributeTok{filename =} \StringTok{"g4\_elasticite\_croisee.png"}\NormalTok{,}
  \AttributeTok{plot =}\NormalTok{ g4,}
  \AttributeTok{width =} \DecValTok{10}\NormalTok{,}
  \AttributeTok{height =} \DecValTok{6}\NormalTok{,}
  \AttributeTok{dpi =} \DecValTok{300}
\NormalTok{)}


\CommentTok{\# sauvegarde du graphique G6}

\FunctionTok{ggsave}\NormalTok{(}
  \AttributeTok{filename =} \StringTok{"g6\_report\_modal\_prix.png"}\NormalTok{,}
  \AttributeTok{plot =}\NormalTok{ g6,}
  \AttributeTok{width =} \DecValTok{10}\NormalTok{,}
  \AttributeTok{height =} \DecValTok{6}\NormalTok{,}
  \AttributeTok{dpi =} \DecValTok{300}
\NormalTok{)}
\end{Highlighting}
\end{Shaded}

Les resultats de la SQ2 sont interessants car ils suggerent l'existence
d'un certain report modal vers les transports en commun, mais avec
plusieurs limites importantes d'interpretation.

Le premier graphique (G4) montre que l'elasticite croisee moyenne est
positive dans les trois zones :

\begin{itemize}
\tightlist
\item
  1,36 en zone urbaine ;
\item
  1,35 en zone periurbaine ;
\item
  1,33 en zone rurale.
\end{itemize}

En theorie economique, une elasticite croisee positive signifie que les
deux biens peuvent etre consideres comme substituts. Dans ce contexte,
cela suggere qu'une hausse du prix du SP95 pourrait etre associee a une
augmentation de l'utilisation des transports en commun.

Cependant, les intervalles de confiance sont tres larges et traversent
zero dans toutes les zones :

\begin{itemize}
\tightlist
\item
  IC ≈ {[}-0,5 ; 3,2{]}.
\end{itemize}

Ce point est important car il indique une forte variabilite statistique.
Ainsi, meme si la moyenne est positive, les resultats ne sont pas
suffisamment stables pour confirmer avec certitude l'existence d'un
report modal fort et systematique.

Le second graphique (G6) est probablement le plus interessant de cette
section.

Il montre que :

\begin{itemize}
\tightlist
\item
  les elasticites croisees sont tres positives dans les classes de prix
  faibles (F1 et F2) ;
\item
  elles deviennent ensuite negatives en F3 et F4 ;
\item
  avant de redevenir legerement positives en F5.
\end{itemize}

Cela suggere que la relation entre le prix du carburant et l'utilisation
des transports en commun n'est pas lineaire.

Plusieurs interpretations peuvent etre proposees.

Lorsque les prix sont relativement faibles, de petites variations du
SP95 peuvent coincider avec des periodes de croissance economique ou
d'augmentation generale de la mobilite. Dans ce contexte, le trafic
routier et les transports en commun peuvent augmenter simultanement.

A partir de F3 (\textasciitilde1,70€--1,90€), les elasticites deviennent
negatives. Cela peut indiquer que la hausse du prix commence a affecter
plus globalement la mobilite, en reduisant a la fois l'utilisation de la
voiture et certains deplacements en transports en commun.

La forte baisse observee en F4 (\textasciitilde1,90€--2,10€) peut
egalement refleter le contexte de crise energetique et d'inflation, dans
lequel l'augmentation du cout de la mobilite affecte plusieurs formes de
deplacement en meme temps.

Enfin, un autre element important est que les trois zones presentent des
resultats presque identiques. Cela suggere soit un comportement agregé
relativement similaire entre les zones, soit une structure du dataset
produisant des dynamiques tres proches entre les differents territoires.

\subsubsection{SQ3 --- facteur geographique (contrainte vs
choix)}\label{sq3-facteur-geographique-contrainte-vs-choix}

l'elasticite est-elle identique en zone rurale et en zone urbaine ?

\begin{Shaded}
\begin{Highlighting}[]
\FunctionTok{cat}\NormalTok{(}\StringTok{"}\SpecialCharTok{\textbackslash{}n}\StringTok{━━━━━━━━━━━━━━━━━━━━━━━━━━━━━━━━━━━━━━━━━━━━━━━━━━━━━━━━━}\SpecialCharTok{\textbackslash{}n}\StringTok{"}\NormalTok{)}
\end{Highlighting}
\end{Shaded}

\begin{verbatim}
## 
## ━━━━━━━━━━━━━━━━━━━━━━━━━━━━━━━━━━━━━━━━━━━━━━━━━━━━━━━━━
\end{verbatim}

\begin{Shaded}
\begin{Highlighting}[]
\FunctionTok{cat}\NormalTok{(}\StringTok{"  SQ3 {-} facteur geographique}\SpecialCharTok{\textbackslash{}n}\StringTok{"}\NormalTok{)}
\end{Highlighting}
\end{Shaded}

\begin{verbatim}
##   SQ3 - facteur geographique
\end{verbatim}

\begin{Shaded}
\begin{Highlighting}[]
\FunctionTok{cat}\NormalTok{(}\StringTok{"  A. elasticite directe discrete par zone}\SpecialCharTok{\textbackslash{}n}\StringTok{"}\NormalTok{)}
\end{Highlighting}
\end{Shaded}

\begin{verbatim}
##   A. elasticite directe discrete par zone
\end{verbatim}

\begin{Shaded}
\begin{Highlighting}[]
\FunctionTok{cat}\NormalTok{(}\StringTok{"  B. regression log{-}lineaire structurelle (nijkamp \& gelman)}\SpecialCharTok{\textbackslash{}n}\StringTok{"}\NormalTok{)}
\end{Highlighting}
\end{Shaded}

\begin{verbatim}
##   B. regression log-lineaire structurelle (nijkamp & gelman)
\end{verbatim}

\begin{Shaded}
\begin{Highlighting}[]
\FunctionTok{cat}\NormalTok{(}\StringTok{"━━━━━━━━━━━━━━━━━━━━━━━━━━━━━━━━━━━━━━━━━━━━━━━━━━━━━━━━━}\SpecialCharTok{\textbackslash{}n}\StringTok{"}\NormalTok{)}
\end{Highlighting}
\end{Shaded}

\begin{verbatim}
## ━━━━━━━━━━━━━━━━━━━━━━━━━━━━━━━━━━━━━━━━━━━━━━━━━━━━━━━━━
\end{verbatim}

\begin{Shaded}
\begin{Highlighting}[]
\CommentTok{\# ── SQ3{-}A : elasticite directe discrete ──────────────────────────────────────}

\NormalTok{elast\_zone }\OtherTok{\textless{}{-}}\NormalTok{ df\_seuil }\SpecialCharTok{\%\textgreater{}\%}

  \FunctionTok{group\_by}\NormalTok{(Zone) }\SpecialCharTok{\%\textgreater{}\%}

  \FunctionTok{summarise}\NormalTok{(}

    \CommentTok{\# nombre d\textquotesingle{}observations par zone}

    \AttributeTok{n =} \FunctionTok{n}\NormalTok{(),}

    \CommentTok{\# elasticite{-}prix moyenne du SP95}

    \AttributeTok{eps\_moy\_SP95 =} \FunctionTok{mean}\NormalTok{(epsilon\_SP95, }\AttributeTok{na.rm =} \ConstantTok{TRUE}\NormalTok{),}

    \CommentTok{\# elasticite{-}prix mediane du SP95}

    \AttributeTok{eps\_med\_SP95 =} \FunctionTok{median}\NormalTok{(epsilon\_SP95, }\AttributeTok{na.rm =} \ConstantTok{TRUE}\NormalTok{),}

    \CommentTok{\# ecart{-}type de l\textquotesingle{}elasticite{-}prix du SP95}

    \AttributeTok{eps\_sd\_SP95 =} \FunctionTok{sd}\NormalTok{(epsilon\_SP95, }\AttributeTok{na.rm =} \ConstantTok{TRUE}\NormalTok{),}

    \CommentTok{\# elasticite moyenne par rapport au gazole}
    \CommentTok{\# ici, on regarde le rapport variation trafic / variation gazole}

    \AttributeTok{eps\_moy\_Gaz =} \FunctionTok{mean}\NormalTok{(Var.Trafic.. }\SpecialCharTok{/}\NormalTok{ Var.Gazole.., }\AttributeTok{na.rm =} \ConstantTok{TRUE}\NormalTok{),}

    \AttributeTok{.groups =} \StringTok{"drop"}
\NormalTok{  ) }\SpecialCharTok{\%\textgreater{}\%}

  \FunctionTok{mutate}\NormalTok{(}

    \CommentTok{\# intervalle de confiance a 95 \%}

    \AttributeTok{IC\_inf =}\NormalTok{ eps\_moy\_SP95 }\SpecialCharTok{{-}} \FloatTok{1.96} \SpecialCharTok{*}\NormalTok{ eps\_sd\_SP95 }\SpecialCharTok{/} \FunctionTok{sqrt}\NormalTok{(n),}
    \AttributeTok{IC\_sup =}\NormalTok{ eps\_moy\_SP95 }\SpecialCharTok{+} \FloatTok{1.96} \SpecialCharTok{*}\NormalTok{ eps\_sd\_SP95 }\SpecialCharTok{/} \FunctionTok{sqrt}\NormalTok{(n),}

    \CommentTok{\# interpretation automatique selon l\textquotesingle{}intensite de l\textquotesingle{}elasticite}

    \AttributeTok{Interpretation =} \FunctionTok{case\_when}\NormalTok{(}
      \FunctionTok{abs}\NormalTok{(eps\_moy\_SP95) }\SpecialCharTok{\textless{}} \FloatTok{0.15} \SpecialCharTok{\textasciitilde{}} \StringTok{"tres inelastique (dependance forte)"}\NormalTok{,}
      \FunctionTok{abs}\NormalTok{(eps\_moy\_SP95) }\SpecialCharTok{\textless{}} \FloatTok{0.40} \SpecialCharTok{\textasciitilde{}} \StringTok{"peu elastique"}\NormalTok{,}
      \FunctionTok{abs}\NormalTok{(eps\_moy\_SP95) }\SpecialCharTok{\textless{}} \FloatTok{1.00} \SpecialCharTok{\textasciitilde{}} \StringTok{"elasticite moderee"}\NormalTok{,}
      \ConstantTok{TRUE} \SpecialCharTok{\textasciitilde{}} \StringTok{"tres elastique (sensibilite forte)"}
\NormalTok{    )}
\NormalTok{  )}


\CommentTok{\# affichage du tableau de synthese}

\FunctionTok{cat}\NormalTok{(}\StringTok{"}\SpecialCharTok{\textbackslash{}n}\StringTok{ tableau de synthese SQ3{-}A {-} elasticite directe par zone :}\SpecialCharTok{\textbackslash{}n}\StringTok{"}\NormalTok{)}
\end{Highlighting}
\end{Shaded}

\begin{verbatim}
## 
##  tableau de synthese SQ3-A - elasticite directe par zone :
\end{verbatim}

\begin{Shaded}
\begin{Highlighting}[]
\NormalTok{elast\_zone }\SpecialCharTok{\%\textgreater{}\%}

  \FunctionTok{select}\NormalTok{(}
\NormalTok{    Zone,}
\NormalTok{    eps\_moy\_SP95,}
\NormalTok{    eps\_med\_SP95,}
\NormalTok{    IC\_inf,}
\NormalTok{    IC\_sup,}
\NormalTok{    eps\_moy\_Gaz,}
\NormalTok{    Interpretation}
\NormalTok{  ) }\SpecialCharTok{\%\textgreater{}\%}

  \FunctionTok{rename}\NormalTok{(}
    \StringTok{"epsilon SP95 (moy)"} \OtherTok{=}\NormalTok{ eps\_moy\_SP95,}
    \StringTok{"epsilon SP95 (med)"} \OtherTok{=}\NormalTok{ eps\_med\_SP95,}
    \StringTok{"IC 95\% inf"} \OtherTok{=}\NormalTok{ IC\_inf,}
    \StringTok{"IC 95\% sup"} \OtherTok{=}\NormalTok{ IC\_sup,}
    \StringTok{"epsilon gazole (moy)"} \OtherTok{=}\NormalTok{ eps\_moy\_Gaz}
\NormalTok{  ) }\SpecialCharTok{\%\textgreater{}\%}

  \FunctionTok{mutate}\NormalTok{(}
    \FunctionTok{across}\NormalTok{(}\FunctionTok{where}\NormalTok{(is.numeric), }\SpecialCharTok{\textasciitilde{}} \FunctionTok{round}\NormalTok{(.x, }\DecValTok{4}\NormalTok{))}
\NormalTok{  ) }\SpecialCharTok{\%\textgreater{}\%}

  \FunctionTok{print}\NormalTok{()}
\end{Highlighting}
\end{Shaded}

\begin{verbatim}
## # A tibble: 3 x 7
##   Zone       `epsilon SP95 (moy)` `epsilon SP95 (med)` `IC 95% inf` `IC 95% sup`
##   <fct>                     <dbl>                <dbl>        <dbl>        <dbl>
## 1 Urbaine                    1.11                0.384       -0.163         2.38
## 2 Périurbai~                 1.06                0.308       -0.187         2.31
## 3 Rurale                     1.02                0.135       -0.212         2.24
## # i 2 more variables: `epsilon gazole (moy)` <dbl>, Interpretation <chr>
\end{verbatim}

\begin{Shaded}
\begin{Highlighting}[]
\CommentTok{\# creation du graphique SQ3{-}A}

\NormalTok{g2 }\OtherTok{\textless{}{-}} \FunctionTok{ggplot}\NormalTok{(elast\_zone, }\FunctionTok{aes}\NormalTok{(}\AttributeTok{x =}\NormalTok{ Zone, }\AttributeTok{y =}\NormalTok{ eps\_moy\_SP95, }\AttributeTok{fill =}\NormalTok{ Zone)) }\SpecialCharTok{+}

  \CommentTok{\# barres des elasticites moyennes}

  \FunctionTok{geom\_col}\NormalTok{(}\AttributeTok{alpha =} \FloatTok{0.85}\NormalTok{, }\AttributeTok{width =} \FloatTok{0.6}\NormalTok{) }\SpecialCharTok{+}

  \CommentTok{\# barres d\textquotesingle{}erreur : intervalle de confiance}

  \FunctionTok{geom\_errorbar}\NormalTok{(}
    \FunctionTok{aes}\NormalTok{(}\AttributeTok{ymin =}\NormalTok{ IC\_inf, }\AttributeTok{ymax =}\NormalTok{ IC\_sup),}
    \AttributeTok{width =} \FloatTok{0.2}\NormalTok{,}
    \AttributeTok{linewidth =} \FloatTok{0.8}
\NormalTok{  ) }\SpecialCharTok{+}

  \CommentTok{\# ligne de reference a zero}

  \FunctionTok{geom\_hline}\NormalTok{(}
    \AttributeTok{yintercept =} \DecValTok{0}\NormalTok{,}
    \AttributeTok{linetype =} \StringTok{"solid"}
\NormalTok{  ) }\SpecialCharTok{+}

  \CommentTok{\# ligne de reference pour l\textquotesingle{}elasticite unitaire}

  \FunctionTok{geom\_hline}\NormalTok{(}
    \AttributeTok{yintercept =} \SpecialCharTok{{-}}\DecValTok{1}\NormalTok{,}
    \AttributeTok{linetype =} \StringTok{"dashed"}\NormalTok{,}
    \AttributeTok{colour =} \StringTok{"\#F44336"}\NormalTok{,}
    \AttributeTok{alpha =} \FloatTok{0.6}
\NormalTok{  ) }\SpecialCharTok{+}

  \CommentTok{\# annotation de l\textquotesingle{}elasticite unitaire}

  \FunctionTok{annotate}\NormalTok{(}
    \StringTok{"text"}\NormalTok{,}
    \AttributeTok{x =} \FloatTok{3.4}\NormalTok{,}
    \AttributeTok{y =} \SpecialCharTok{{-}}\FloatTok{1.05}\NormalTok{,}
    \AttributeTok{label =} \StringTok{"elasticite unitaire"}\NormalTok{,}
    \AttributeTok{size =} \DecValTok{3}\NormalTok{,}
    \AttributeTok{colour =} \StringTok{"\#F44336"}
\NormalTok{  ) }\SpecialCharTok{+}

  \CommentTok{\# couleurs par zone}

  \FunctionTok{scale\_fill\_manual}\NormalTok{(}\AttributeTok{values =}\NormalTok{ couleurs\_zone) }\SpecialCharTok{+}

  \CommentTok{\# titres et axes}

  \FunctionTok{labs}\NormalTok{(}
    \AttributeTok{title =} \StringTok{"G2 {-} SQ3 : elasticite{-}prix directe de la demande de trafic (SP95)"}\NormalTok{,}
    \AttributeTok{subtitle =} \StringTok{"epsilon = delta trafic \% / delta prix \%  |  IC 95\%  |  rurale \textasciitilde{} 0 {-}\textgreater{} contrainte"}\NormalTok{,}
    \AttributeTok{x =} \ConstantTok{NULL}\NormalTok{,}
    \AttributeTok{y =} \StringTok{"Elasticite (epsilon)"}\NormalTok{,}
    \AttributeTok{fill =} \ConstantTok{NULL}
\NormalTok{  ) }\SpecialCharTok{+}

  \CommentTok{\# mise en forme generale}

  \FunctionTok{theme\_minimal}\NormalTok{(}\AttributeTok{base\_size =} \DecValTok{12}\NormalTok{) }\SpecialCharTok{+}

  \FunctionTok{theme}\NormalTok{(}
    \AttributeTok{legend.position =} \StringTok{"none"}
\NormalTok{  )}


\CommentTok{\# affichage du graphique}

\NormalTok{g2}
\end{Highlighting}
\end{Shaded}

\pandocbounded{\includegraphics[keepaspectratio]{price_elasticity_gasoline_demand_corrigé_files/figure-latex/unnamed-chunk-21-1.pdf}}

\begin{Shaded}
\begin{Highlighting}[]
\CommentTok{\# sauvegarde du graphique}

\FunctionTok{ggsave}\NormalTok{(}
  \AttributeTok{filename =} \StringTok{"g2\_sq3\_elasticite\_directe.png"}\NormalTok{,}
  \AttributeTok{plot =}\NormalTok{ g2,}
  \AttributeTok{width =} \DecValTok{10}\NormalTok{,}
  \AttributeTok{height =} \DecValTok{6}\NormalTok{,}
  \AttributeTok{dpi =} \DecValTok{300}
\NormalTok{)}
\end{Highlighting}
\end{Shaded}

Ces resultats montrent, de maniere generale, qu'il n'existe pas de
difference geographique tres forte entre les zones, et que l'elasticite
apparait positive dans les trois cas, mais avec une incertitude
statistique importante.

Le premier element marquant est la proximite des moyennes :

Urbaine : 1,11 ; Periurbaine : 1,06 ; Rurale : 1,02.

Cela suggere que le comportement observe est relativement similaire dans
les trois zones.

Les intervalles de confiance sont cependant tres larges et traversent
tous zero :

Urbaine : {[}-0,163 ; 2,38{]} ; Periurbaine : {[}-0,187 ; 2,31{]} ;
Rurale : {[}-0,212 ; 2,24{]}.

D'un point de vue statistique, cela signifie que les estimations restent
fragiles. Il n'est donc pas possible d'affirmer avec certitude que
l'elasticite est reellement positive et stable.

Par ailleurs, les medianes sont nettement plus faibles que les moyennes
:

0,384 ; 0,308 ; 0,135.

Cela indique que la distribution contient quelques valeurs elevees qui
tirent artificiellement la moyenne vers le haut. En d'autres termes, la
moyenne donne l'impression d'une elasticite plus forte que celle
observee dans la majorite des cas.

Le graphique confirme cette interpretation : les barres se situent
autour de 1, mais les intervalles de confiance sont tres importants, ce
qui traduit une forte variabilite des observations.

Ainsi, l'interpretation la plus prudente consiste a dire que le
comportement geographique ne semble pas tres different entre les zones,
et que l'elasticite directe du trafic par rapport au prix du SP95 reste
instable, fortement dispersee et statistiquement peu robuste.

Enfin, un point important doit etre souligne : comme l'elasticite
moyenne est positive, le resultat ne correspond pas au comportement
economique classique d'une demande qui diminue lorsque le prix augmente.
Cela renforce l'idee que les estimations sont probablement influencees
par des tendances temporelles, des chocs externes ou des variables
omises, et qu'elles ne doivent pas etre interpretees comme une relation
causale pure entre le prix du carburant et le trafic routier.

\subsubsection{SQ3-B --- regression log-lineaire
structurelle}\label{sq3-b-regression-log-lineaire-structurelle}

l'elasticite structurelle du trafic varie-t-elle selon la zone ?

\begin{Shaded}
\begin{Highlighting}[]
\CommentTok{\# ── harmonisation des noms de colonnes pour SQ3{-}B ────────────────────────────}
\CommentTok{\# on cree un dataframe temporaire avec des noms standardises}
\CommentTok{\# cela evite les erreurs si le fichier contient encore les anciens noms}

\NormalTok{df\_abs\_sq3 }\OtherTok{\textless{}{-}}\NormalTok{ df\_abs}

\CommentTok{\# trafic routier}
\NormalTok{df\_abs\_sq3}\SpecialCharTok{$}\NormalTok{Trafic }\OtherTok{\textless{}{-}} \ControlFlowTok{if}\NormalTok{ (}\StringTok{"Trafic"} \SpecialCharTok{\%in\%} \FunctionTok{names}\NormalTok{(df\_abs\_sq3)) \{}
\NormalTok{  df\_abs\_sq3}\SpecialCharTok{$}\NormalTok{Trafic}
\NormalTok{\} }\ControlFlowTok{else} \ControlFlowTok{if}\NormalTok{ (}\StringTok{"Trafic.Routier"} \SpecialCharTok{\%in\%} \FunctionTok{names}\NormalTok{(df\_abs\_sq3)) \{}
\NormalTok{  df\_abs\_sq3}\SpecialCharTok{$}\NormalTok{Trafic.Routier}
\NormalTok{\} }\ControlFlowTok{else}\NormalTok{ \{}
  \FunctionTok{stop}\NormalTok{(}\StringTok{"colonne du trafic introuvable dans df\_abs"}\NormalTok{)}
\NormalTok{\}}

\CommentTok{\# prix du gazole}
\NormalTok{df\_abs\_sq3}\SpecialCharTok{$}\NormalTok{Prix\_Gazole }\OtherTok{\textless{}{-}} \ControlFlowTok{if}\NormalTok{ (}\StringTok{"Prix\_Gazole"} \SpecialCharTok{\%in\%} \FunctionTok{names}\NormalTok{(df\_abs\_sq3)) \{}
\NormalTok{  df\_abs\_sq3}\SpecialCharTok{$}\NormalTok{Prix\_Gazole}
\NormalTok{\} }\ControlFlowTok{else} \ControlFlowTok{if}\NormalTok{ (}\StringTok{"Prix.Gazole"} \SpecialCharTok{\%in\%} \FunctionTok{names}\NormalTok{(df\_abs\_sq3)) \{}
\NormalTok{  df\_abs\_sq3}\SpecialCharTok{$}\NormalTok{Prix.Gazole}
\NormalTok{\} }\ControlFlowTok{else}\NormalTok{ \{}
  \FunctionTok{stop}\NormalTok{(}\StringTok{"colonne du prix du gazole introuvable dans df\_abs"}\NormalTok{)}
\NormalTok{\}}

\CommentTok{\# voyageurs en transports en commun}
\NormalTok{df\_abs\_sq3}\SpecialCharTok{$}\NormalTok{Voyageurs }\OtherTok{\textless{}{-}} \ControlFlowTok{if}\NormalTok{ (}\StringTok{"Voyageurs"} \SpecialCharTok{\%in\%} \FunctionTok{names}\NormalTok{(df\_abs\_sq3)) \{}
\NormalTok{  df\_abs\_sq3}\SpecialCharTok{$}\NormalTok{Voyageurs}
\NormalTok{\} }\ControlFlowTok{else} \ControlFlowTok{if}\NormalTok{ (}\StringTok{"Voyageurs.TC"} \SpecialCharTok{\%in\%} \FunctionTok{names}\NormalTok{(df\_abs\_sq3)) \{}
\NormalTok{  df\_abs\_sq3}\SpecialCharTok{$}\NormalTok{Voyageurs.TC}
\NormalTok{\} }\ControlFlowTok{else}\NormalTok{ \{}
  \FunctionTok{stop}\NormalTok{(}\StringTok{"colonne des voyageurs TC introuvable dans df\_abs"}\NormalTok{)}
\NormalTok{\}}


\CommentTok{\# ── SQ3{-}B : regression log{-}lineaire structurelle ─────────────────────────────}
\CommentTok{\# ln(Trafic) = alpha + beta * ln(Prix\_SP95)}
\CommentTok{\# beta = elasticite structurelle du trafic}

\FunctionTok{cat}\NormalTok{(}\StringTok{"}\SpecialCharTok{\textbackslash{}n}\StringTok{━━━━━━━━━━━━━━━━━━━━━━━━━━━━━━━━━━━━━━━━━━━━━━━━━━━━━━━━━}\SpecialCharTok{\textbackslash{}n}\StringTok{"}\NormalTok{)}
\end{Highlighting}
\end{Shaded}

\begin{verbatim}
## 
## ━━━━━━━━━━━━━━━━━━━━━━━━━━━━━━━━━━━━━━━━━━━━━━━━━━━━━━━━━
\end{verbatim}

\begin{Shaded}
\begin{Highlighting}[]
\FunctionTok{cat}\NormalTok{(}\StringTok{"  SQ3{-}B {-} regression log{-}lineaire structurelle}\SpecialCharTok{\textbackslash{}n}\StringTok{"}\NormalTok{)}
\end{Highlighting}
\end{Shaded}

\begin{verbatim}
##   SQ3-B - regression log-lineaire structurelle
\end{verbatim}

\begin{Shaded}
\begin{Highlighting}[]
\FunctionTok{cat}\NormalTok{(}\StringTok{"  ln(Trafic) = alpha + beta * ln(Prix\_SP95)}\SpecialCharTok{\textbackslash{}n}\StringTok{"}\NormalTok{)}
\end{Highlighting}
\end{Shaded}

\begin{verbatim}
##   ln(Trafic) = alpha + beta * ln(Prix_SP95)
\end{verbatim}

\begin{Shaded}
\begin{Highlighting}[]
\FunctionTok{cat}\NormalTok{(}\StringTok{"━━━━━━━━━━━━━━━━━━━━━━━━━━━━━━━━━━━━━━━━━━━━━━━━━━━━━━━━━}\SpecialCharTok{\textbackslash{}n}\StringTok{"}\NormalTok{)}
\end{Highlighting}
\end{Shaded}

\begin{verbatim}
## ━━━━━━━━━━━━━━━━━━━━━━━━━━━━━━━━━━━━━━━━━━━━━━━━━━━━━━━━━
\end{verbatim}

\begin{Shaded}
\begin{Highlighting}[]
\CommentTok{\# ── estimation des modeles par zone ──────────────────────────────────────────}
\CommentTok{\# 1) trafic \textasciitilde{} prix SP95}
\CommentTok{\# 2) trafic \textasciitilde{} prix gazole}
\CommentTok{\# 3) voyageurs TC \textasciitilde{} prix SP95}

\NormalTok{modeles\_loglin }\OtherTok{\textless{}{-}}\NormalTok{ df\_abs\_sq3 }\SpecialCharTok{\%\textgreater{}\%}
  \FunctionTok{group\_by}\NormalTok{(Zone) }\SpecialCharTok{\%\textgreater{}\%}
  \FunctionTok{nest}\NormalTok{() }\SpecialCharTok{\%\textgreater{}\%}
  \FunctionTok{mutate}\NormalTok{(}

    \CommentTok{\# modele principal : prix SP95 {-}\textgreater{} trafic}
    \AttributeTok{modele\_SP95 =} \FunctionTok{map}\NormalTok{(}
\NormalTok{      data,}
      \SpecialCharTok{\textasciitilde{}} \FunctionTok{lm}\NormalTok{(}\FunctionTok{log}\NormalTok{(Trafic) }\SpecialCharTok{\textasciitilde{}} \FunctionTok{log}\NormalTok{(Prix\_SP95), }\AttributeTok{data =}\NormalTok{ .x)}
\NormalTok{    ),}

    \CommentTok{\# modele complementaire : prix gazole {-}\textgreater{} trafic}
    \AttributeTok{modele\_Gazole =} \FunctionTok{map}\NormalTok{(}
\NormalTok{      data,}
      \SpecialCharTok{\textasciitilde{}} \FunctionTok{lm}\NormalTok{(}\FunctionTok{log}\NormalTok{(Trafic) }\SpecialCharTok{\textasciitilde{}} \FunctionTok{log}\NormalTok{(Prix\_Gazole), }\AttributeTok{data =}\NormalTok{ .x)}
\NormalTok{    ),}

    \CommentTok{\# modele transports en commun : prix SP95 {-}\textgreater{} voyageurs}
    \AttributeTok{modele\_TC =} \FunctionTok{map}\NormalTok{(}
\NormalTok{      data,}
      \SpecialCharTok{\textasciitilde{}} \FunctionTok{lm}\NormalTok{(}\FunctionTok{log}\NormalTok{(Voyageurs }\SpecialCharTok{+} \DecValTok{1}\NormalTok{) }\SpecialCharTok{\textasciitilde{}} \FunctionTok{log}\NormalTok{(Prix\_SP95), }\AttributeTok{data =}\NormalTok{ .x)}
\NormalTok{    ),}

    \CommentTok{\# extraction des coefficients et des statistiques}
    \AttributeTok{tidied\_SP95  =} \FunctionTok{map}\NormalTok{(modele\_SP95, tidy, }\AttributeTok{conf.int =} \ConstantTok{TRUE}\NormalTok{),}
    \AttributeTok{glanced\_SP95 =} \FunctionTok{map}\NormalTok{(modele\_SP95, glance),}

    \AttributeTok{tidied\_Gazole  =} \FunctionTok{map}\NormalTok{(modele\_Gazole, tidy, }\AttributeTok{conf.int =} \ConstantTok{TRUE}\NormalTok{),}
    \AttributeTok{glanced\_Gazole =} \FunctionTok{map}\NormalTok{(modele\_Gazole, glance),}

    \AttributeTok{tidied\_TC  =} \FunctionTok{map}\NormalTok{(modele\_TC, tidy, }\AttributeTok{conf.int =} \ConstantTok{TRUE}\NormalTok{),}
    \AttributeTok{glanced\_TC =} \FunctionTok{map}\NormalTok{(modele\_TC, glance)}
\NormalTok{  )}


\CommentTok{\# ── resultats du modele SP95 {-}\textgreater{} trafic ───────────────────────────────────────}

\NormalTok{res\_loglin\_SP95 }\OtherTok{\textless{}{-}}\NormalTok{ modeles\_loglin }\SpecialCharTok{\%\textgreater{}\%}
  \FunctionTok{unnest}\NormalTok{(tidied\_SP95) }\SpecialCharTok{\%\textgreater{}\%}
  \FunctionTok{filter}\NormalTok{(term }\SpecialCharTok{==} \StringTok{"log(Prix\_SP95)"}\NormalTok{) }\SpecialCharTok{\%\textgreater{}\%}
  \FunctionTok{left\_join}\NormalTok{(}
\NormalTok{    modeles\_loglin }\SpecialCharTok{\%\textgreater{}\%}
      \FunctionTok{unnest}\NormalTok{(glanced\_SP95) }\SpecialCharTok{\%\textgreater{}\%}
      \FunctionTok{select}\NormalTok{(Zone, r.squared, adj.r.squared),}
    \AttributeTok{by =} \StringTok{"Zone"}
\NormalTok{  ) }\SpecialCharTok{\%\textgreater{}\%}
  \FunctionTok{select}\NormalTok{(}
\NormalTok{    Zone, estimate, std.error, statistic, p.value,}
\NormalTok{    conf.low, conf.high, r.squared, adj.r.squared}
\NormalTok{  )}


\CommentTok{\# ── resultats du modele SP95 {-}\textgreater{} TC ───────────────────────────────────────────}

\NormalTok{res\_loglin\_TC }\OtherTok{\textless{}{-}}\NormalTok{ modeles\_loglin }\SpecialCharTok{\%\textgreater{}\%}
  \FunctionTok{unnest}\NormalTok{(tidied\_TC) }\SpecialCharTok{\%\textgreater{}\%}
  \FunctionTok{filter}\NormalTok{(term }\SpecialCharTok{==} \StringTok{"log(Prix\_SP95)"}\NormalTok{) }\SpecialCharTok{\%\textgreater{}\%}
  \FunctionTok{left\_join}\NormalTok{(}
\NormalTok{    modeles\_loglin }\SpecialCharTok{\%\textgreater{}\%}
      \FunctionTok{unnest}\NormalTok{(glanced\_TC) }\SpecialCharTok{\%\textgreater{}\%}
      \FunctionTok{select}\NormalTok{(Zone, r.squared),}
    \AttributeTok{by =} \StringTok{"Zone"}
\NormalTok{  ) }\SpecialCharTok{\%\textgreater{}\%}
  \FunctionTok{select}\NormalTok{(}
\NormalTok{    Zone, estimate, std.error, p.value,}
\NormalTok{    conf.low, conf.high, r.squared}
\NormalTok{  )}


\CommentTok{\# ── affichage des resultats SP95 {-}\textgreater{} trafic ───────────────────────────────────}

\FunctionTok{cat}\NormalTok{(}\StringTok{"}\SpecialCharTok{\textbackslash{}n}\StringTok{SQ3{-}B {-} regression log{-}lineaire (SP95 {-}\textgreater{} trafic)}\SpecialCharTok{\textbackslash{}n}\StringTok{"}\NormalTok{)}
\end{Highlighting}
\end{Shaded}

\begin{verbatim}
## 
## SQ3-B - regression log-lineaire (SP95 -> trafic)
\end{verbatim}

\begin{Shaded}
\begin{Highlighting}[]
\NormalTok{res\_loglin\_SP95 }\SpecialCharTok{\%\textgreater{}\%}
  \FunctionTok{mutate}\NormalTok{(}
    \AttributeTok{Significatif =} \FunctionTok{case\_when}\NormalTok{(}
\NormalTok{      p.value }\SpecialCharTok{\textless{}} \FloatTok{0.001} \SpecialCharTok{\textasciitilde{}} \StringTok{"***"}\NormalTok{,}
\NormalTok{      p.value }\SpecialCharTok{\textless{}} \FloatTok{0.01}  \SpecialCharTok{\textasciitilde{}} \StringTok{"**"}\NormalTok{,}
\NormalTok{      p.value }\SpecialCharTok{\textless{}} \FloatTok{0.05}  \SpecialCharTok{\textasciitilde{}} \StringTok{"*"}\NormalTok{,}
      \ConstantTok{TRUE} \SpecialCharTok{\textasciitilde{}} \StringTok{"ns"}
\NormalTok{    ),}
    \AttributeTok{Interpretation =} \FunctionTok{case\_when}\NormalTok{(}
      \FunctionTok{abs}\NormalTok{(estimate) }\SpecialCharTok{\textless{}} \FloatTok{0.15} \SpecialCharTok{\textasciitilde{}} \StringTok{"tres inelastique"}\NormalTok{,}
      \FunctionTok{abs}\NormalTok{(estimate) }\SpecialCharTok{\textless{}} \FloatTok{0.40} \SpecialCharTok{\textasciitilde{}} \StringTok{"peu elastique"}\NormalTok{,}
      \FunctionTok{abs}\NormalTok{(estimate) }\SpecialCharTok{\textless{}} \FloatTok{1.00} \SpecialCharTok{\textasciitilde{}} \StringTok{"elasticite moderee"}\NormalTok{,}
      \ConstantTok{TRUE} \SpecialCharTok{\textasciitilde{}} \StringTok{"tres elastique"}
\NormalTok{    )}
\NormalTok{  ) }\SpecialCharTok{\%\textgreater{}\%}
  \FunctionTok{mutate}\NormalTok{(}\FunctionTok{across}\NormalTok{(}\FunctionTok{where}\NormalTok{(is.numeric), }\SpecialCharTok{\textasciitilde{}} \FunctionTok{round}\NormalTok{(.x, }\DecValTok{4}\NormalTok{))) }\SpecialCharTok{\%\textgreater{}\%}
  \FunctionTok{print}\NormalTok{()}
\end{Highlighting}
\end{Shaded}

\begin{verbatim}
## # A tibble: 3 x 11
## # Groups:   Zone [3]
##   Zone        estimate std.error statistic p.value conf.low conf.high r.squared
##   <fct>          <dbl>     <dbl>     <dbl>   <dbl>    <dbl>     <dbl>     <dbl>
## 1 Urbaine         1.39     0.222      6.23       0    0.944      1.83     0.314
## 2 Périurbaine     1.23     0.221      5.56       0    0.788      1.67     0.266
## 3 Rurale          1.07     0.220      4.85       0    0.629      1.50     0.217
## # i 3 more variables: adj.r.squared <dbl>, Significatif <chr>,
## #   Interpretation <chr>
\end{verbatim}

\begin{Shaded}
\begin{Highlighting}[]
\CommentTok{\# ── affichage des resultats SP95 {-}\textgreater{} TC ───────────────────────────────────────}

\FunctionTok{cat}\NormalTok{(}\StringTok{"}\SpecialCharTok{\textbackslash{}n}\StringTok{SQ3{-}B {-} regression log{-}lineaire (SP95 {-}\textgreater{} TC)}\SpecialCharTok{\textbackslash{}n}\StringTok{"}\NormalTok{)}
\end{Highlighting}
\end{Shaded}

\begin{verbatim}
## 
## SQ3-B - regression log-lineaire (SP95 -> TC)
\end{verbatim}

\begin{Shaded}
\begin{Highlighting}[]
\NormalTok{res\_loglin\_TC }\SpecialCharTok{\%\textgreater{}\%}
  \FunctionTok{mutate}\NormalTok{(}\FunctionTok{across}\NormalTok{(}\FunctionTok{where}\NormalTok{(is.numeric), }\SpecialCharTok{\textasciitilde{}} \FunctionTok{round}\NormalTok{(.x, }\DecValTok{4}\NormalTok{))) }\SpecialCharTok{\%\textgreater{}\%}
  \FunctionTok{print}\NormalTok{()}
\end{Highlighting}
\end{Shaded}

\begin{verbatim}
## # A tibble: 3 x 7
## # Groups:   Zone [3]
##   Zone        estimate std.error p.value conf.low conf.high r.squared
##   <fct>          <dbl>     <dbl>   <dbl>    <dbl>     <dbl>     <dbl>
## 1 Urbaine         2.22     0.295       0     1.63      2.81     0.400
## 2 Périurbaine     2.22     0.294       0     1.63      2.80     0.4  
## 3 Rurale          2.19     0.290       0     1.62      2.77     0.402
\end{verbatim}

\begin{Shaded}
\begin{Highlighting}[]
\CommentTok{\# ── graphique principal de la regression log{-}lineaire ────────────────────────}

\NormalTok{g5 }\OtherTok{\textless{}{-}} \FunctionTok{ggplot}\NormalTok{(}
\NormalTok{  df\_abs\_sq3,}
  \FunctionTok{aes}\NormalTok{(}
    \AttributeTok{x =} \FunctionTok{log}\NormalTok{(Prix\_SP95),}
    \AttributeTok{y =} \FunctionTok{log}\NormalTok{(Trafic),}
    \AttributeTok{colour =}\NormalTok{ Zone}
\NormalTok{  )}
\NormalTok{) }\SpecialCharTok{+}

  \CommentTok{\# nuage de points par zone}
  \FunctionTok{geom\_point}\NormalTok{(}\AttributeTok{alpha =} \FloatTok{0.4}\NormalTok{, }\AttributeTok{size =} \FloatTok{1.5}\NormalTok{) }\SpecialCharTok{+}

  \CommentTok{\# droite de regression lineaire}
  \FunctionTok{geom\_smooth}\NormalTok{(}
    \AttributeTok{method =} \StringTok{"lm"}\NormalTok{,}
    \AttributeTok{se =} \ConstantTok{TRUE}\NormalTok{,}
    \AttributeTok{linewidth =} \FloatTok{1.2}
\NormalTok{  ) }\SpecialCharTok{+}

  \CommentTok{\# couleurs des zones}
  \FunctionTok{scale\_colour\_manual}\NormalTok{(}\AttributeTok{values =}\NormalTok{ couleurs\_zone) }\SpecialCharTok{+}

  \CommentTok{\# equation de regression sur le graphique (ggpmisc)}
  \FunctionTok{stat\_poly\_eq}\NormalTok{(}
    \FunctionTok{aes}\NormalTok{(}\AttributeTok{label =} \FunctionTok{after\_stat}\NormalTok{(eq.label)),}
    \AttributeTok{formula =}\NormalTok{ y }\SpecialCharTok{\textasciitilde{}}\NormalTok{ x,}
    \AttributeTok{parse =} \ConstantTok{TRUE}\NormalTok{,}
    \AttributeTok{size =} \DecValTok{3}\NormalTok{,}
    \AttributeTok{label.x =} \StringTok{"left"}\NormalTok{,}
    \AttributeTok{label.y =} \StringTok{"top"}
\NormalTok{  ) }\SpecialCharTok{+}

  \CommentTok{\# titres et axes}
  \FunctionTok{labs}\NormalTok{(}
    \AttributeTok{title =} \StringTok{"G5 {-} SQ3 : regression log{-}lineaire : ln(Trafic) \textasciitilde{} ln(Prix\_SP95)"}\NormalTok{,}
    \AttributeTok{subtitle =} \StringTok{"la pente beta mesure l\textquotesingle{}elasticite structurelle du trafic"}\NormalTok{,}
    \AttributeTok{x =} \StringTok{"ln(Prix SP95)"}\NormalTok{,}
    \AttributeTok{y =} \StringTok{"ln(Trafic routier)"}\NormalTok{,}
    \AttributeTok{colour =} \StringTok{"Zone"}
\NormalTok{  ) }\SpecialCharTok{+}

  \CommentTok{\# mise en forme generale}
  \FunctionTok{theme\_minimal}\NormalTok{(}\AttributeTok{base\_size =} \DecValTok{12}\NormalTok{) }\SpecialCharTok{+}

  \FunctionTok{theme}\NormalTok{(}
    \AttributeTok{legend.position =} \StringTok{"top"}
\NormalTok{  )}


\CommentTok{\# ── affichage du graphique ───────────────────────────────────────────────────}

\NormalTok{g5}
\end{Highlighting}
\end{Shaded}

\pandocbounded{\includegraphics[keepaspectratio]{price_elasticity_gasoline_demand_corrigé_files/figure-latex/unnamed-chunk-22-1.pdf}}

\begin{Shaded}
\begin{Highlighting}[]
\CommentTok{\# ── sauvegarde du graphique ──────────────────────────────────────────────────}

\FunctionTok{ggsave}\NormalTok{(}
  \AttributeTok{filename =} \StringTok{"g5\_sq3\_regression\_loglineaire.png"}\NormalTok{,}
  \AttributeTok{plot =}\NormalTok{ g5,}
  \AttributeTok{width =} \DecValTok{10}\NormalTok{,}
  \AttributeTok{height =} \DecValTok{6}\NormalTok{,}
  \AttributeTok{dpi =} \DecValTok{300}
\NormalTok{)}
\end{Highlighting}
\end{Shaded}

ce que montre la regression log-lineaire

Le modele estime la relation suivante :

ln(Trafic)=α+β×ln(Prix\_SP95)+ε

Dans ce type de modele, le coefficient beta correspond directement a une
elasticite structurelle constante.

Autrement dit :

si le prix augmente de 1 \%, alors le trafic varie en moyenne de beta
\%. 1. les coefficients beta sont positifs dans toutes les zones

Les elasticites estimees sont :

zone urbaine : 1,39 zone periurbaine : 1,23 zone rurale : 1,07

Cela signifie que, dans le modele estime, une hausse du prix du SP95 est
associee a une hausse du trafic routier.

Ce resultat est contre-intuitif du point de vue economique classique,
car on attend normalement une relation negative entre prix et demande.

Donc :

le modele ne doit pas etre interprete comme une causalite directe
simple.

Il capture plutot une dynamique temporelle globale du dataset :

reprise economique post-covid ; retour progressif de la mobilite ;
inflation energetique ; augmentation simultanee du trafic et des prix
entre 2021 et 2024.

Autrement dit :

le prix et le trafic augmentent ensemble dans le temps, ce qui produit
une pente positive.

\begin{enumerate}
\def\labelenumi{\arabic{enumi}.}
\setcounter{enumi}{1}
\tightlist
\item
  les resultats sont statistiquement significatifs
\end{enumerate}

C'est le point le plus important.

Contrairement aux elasticites discretes precedentes :

les p-values sont quasiment nulles ; les coefficients sont tres
significatifs (***) ; les intervalles de confiance ne traversent pas
zero.

Par exemple :

urbaine : IC {[}0,944 ; 1,83{]} periurbaine : IC {[}0,788 ; 1,67{]}
rurale : IC {[}0,629 ; 1,50{]}

Donc :

la relation positive observee est statistiquement robuste dans le cadre
du modele.

\begin{enumerate}
\def\labelenumi{\arabic{enumi}.}
\setcounter{enumi}{2}
\tightlist
\item
  le facteur geographique reste relativement faible
\end{enumerate}

Les trois coefficients restent proches :

1,39 1,23 1,07

Cela suggere encore une fois que :

les comportements agreges sont assez similaires entre zones.

La zone urbaine apparait legerement plus sensible au prix, tandis que la
zone rurale semble un peu plus stable, mais les differences restent
moderees.

Donc :

le facteur geographique existe, mais il n'introduit pas une rupture
majeure dans les comportements observes.

\begin{enumerate}
\def\labelenumi{\arabic{enumi}.}
\setcounter{enumi}{3}
\tightlist
\item
  le pouvoir explicatif reste moyen
\end{enumerate}

Les R² sont :

0,314 (urbaine) 0,266 (periurbaine) 0,217 (rurale)

Donc le prix explique seulement environ :

22 \% a 31 \% de la variation du trafic.

Cela veut dire que :

une grande partie des variations du trafic depend probablement d'autres
variables :

contexte sanitaire ; saisonnalite ; activite economique ; politiques
publiques ; teletravail ; habitudes de mobilite.

Le prix du carburant n'est donc pas le seul determinant du trafic.

\begin{enumerate}
\def\labelenumi{\arabic{enumi}.}
\setcounter{enumi}{4}
\tightlist
\item
  regression SP95 -\textgreater{} transports en commun
\end{enumerate}

Les coefficients deviennent encore plus eleves :

environ 2,2 dans les trois zones.

Cela suggere qu'une hausse du prix du carburant est associee a une
hausse importante de l'utilisation des transports en commun.

Et ici encore :

les coefficients sont tres significatifs ; les IC restent entierement
positifs ; les R² (\textasciitilde0,40) sont meilleurs que pour le
trafic routier.

Donc :

la relation statistique entre prix du carburant et usage des transports
en commun apparait plus forte que celle observee avec le trafic routier.

lecture globale prudente

La conclusion la plus prudente serait probablement :

les regressions montrent une relation statistique robuste entre hausse
des prix et evolution des mobilites ; cependant, les coefficients
positifs indiquent surtout une dynamique temporelle commune plutot
qu'une elasticite causale pure ; les resultats suggerent aussi que les
transports en commun reagissent davantage aux variations du prix du
carburant que le trafic routier lui-meme.

\subsubsection{grand finale}\label{grand-finale}

\begin{Shaded}
\begin{Highlighting}[]
\CommentTok{\# ══════════════════════════════════════════════════════════════════════════════}
\CommentTok{\#  TABLEAU DE SYNTHÈSE FINAL}
\CommentTok{\# ══════════════════════════════════════════════════════════════════════════════}
\FunctionTok{cat}\NormalTok{(}\StringTok{"}\SpecialCharTok{\textbackslash{}n}\StringTok{━━━━━━━━━━━━━━━━━━━━━━━━━━━━━━━━━━━━━━━━━━━━━━━━━━━━━━━━━}\SpecialCharTok{\textbackslash{}n}\StringTok{"}\NormalTok{)}
\end{Highlighting}
\end{Shaded}

\begin{verbatim}
## 
## ━━━━━━━━━━━━━━━━━━━━━━━━━━━━━━━━━━━━━━━━━━━━━━━━━━━━━━━━━
\end{verbatim}

\begin{Shaded}
\begin{Highlighting}[]
\FunctionTok{cat}\NormalTok{(}\StringTok{"  TABLEAU DE SYNTHESE FINAL (pour le rapport)}\SpecialCharTok{\textbackslash{}n}\StringTok{"}\NormalTok{)}
\end{Highlighting}
\end{Shaded}

\begin{verbatim}
##   TABLEAU DE SYNTHESE FINAL (pour le rapport)
\end{verbatim}

\begin{Shaded}
\begin{Highlighting}[]
\FunctionTok{cat}\NormalTok{(}\StringTok{"━━━━━━━━━━━━━━━━━━━━━━━━━━━━━━━━━━━━━━━━━━━━━━━━━━━━━━━━━}\SpecialCharTok{\textbackslash{}n}\StringTok{"}\NormalTok{)}
\end{Highlighting}
\end{Shaded}

\begin{verbatim}
## ━━━━━━━━━━━━━━━━━━━━━━━━━━━━━━━━━━━━━━━━━━━━━━━━━━━━━━━━━
\end{verbatim}

\begin{Shaded}
\begin{Highlighting}[]
\NormalTok{corr\_zone }\OtherTok{\textless{}{-}}\NormalTok{ df\_abs }\SpecialCharTok{\%\textgreater{}\%}
  \FunctionTok{group\_by}\NormalTok{(Zone) }\SpecialCharTok{\%\textgreater{}\%}
  \FunctionTok{summarise}\NormalTok{(}
    \AttributeTok{corr\_SP95\_trafic =} \FunctionTok{cor}\NormalTok{(Prix\_SP95, Trafic.Routier, }\AttributeTok{use =} \StringTok{"complete.obs"}\NormalTok{),}
    \AttributeTok{.groups =} \StringTok{"drop"}
\NormalTok{  )}
 
\NormalTok{synthese\_finale }\OtherTok{\textless{}{-}}\NormalTok{ corr\_zone }\SpecialCharTok{\%\textgreater{}\%}
  \FunctionTok{left\_join}\NormalTok{(elast\_zone     }\SpecialCharTok{\%\textgreater{}\%} \FunctionTok{select}\NormalTok{(Zone, eps\_moy\_SP95, Interpretation),           }\AttributeTok{by =} \StringTok{"Zone"}\NormalTok{) }\SpecialCharTok{\%\textgreater{}\%}
  \FunctionTok{left\_join}\NormalTok{(cross\_zone     }\SpecialCharTok{\%\textgreater{}\%} \FunctionTok{select}\NormalTok{(Zone, eps\_crois\_moy),                           }\AttributeTok{by =} \StringTok{"Zone"}\NormalTok{) }\SpecialCharTok{\%\textgreater{}\%}
  \FunctionTok{left\_join}\NormalTok{(seuil\_rupture  }\SpecialCharTok{\%\textgreater{}\%} \FunctionTok{select}\NormalTok{(Zone, Faix\_Rupture, Delta\_max),                 }\AttributeTok{by =} \StringTok{"Zone"}\NormalTok{) }\SpecialCharTok{\%\textgreater{}\%}
  \FunctionTok{left\_join}\NormalTok{(res\_loglin\_SP95 }\SpecialCharTok{\%\textgreater{}\%} \FunctionTok{select}\NormalTok{(Zone, }\AttributeTok{beta\_loglin =}\NormalTok{ estimate,}
                                        \AttributeTok{R2 =}\NormalTok{ r.squared, }\AttributeTok{Sig =}\NormalTok{ p.value),             }\AttributeTok{by =} \StringTok{"Zone"}\NormalTok{) }\SpecialCharTok{\%\textgreater{}\%}
  \FunctionTok{rename}\NormalTok{(}
    \StringTok{"Correlation (r)"}    \OtherTok{=}\NormalTok{ corr\_SP95\_trafic,}
    \StringTok{"epsilon directe"}    \OtherTok{=}\NormalTok{ eps\_moy\_SP95,}
    \StringTok{"Interpretation"}     \OtherTok{=}\NormalTok{ Interpretation,}
    \StringTok{"epsilon croisee TC"} \OtherTok{=}\NormalTok{ eps\_crois\_moy,}
    \StringTok{"Seuil rupture"}      \OtherTok{=}\NormalTok{ Faix\_Rupture,}
    \StringTok{"Delta epsilon"}      \OtherTok{=}\NormalTok{ Delta\_max,}
    \StringTok{"beta log{-}lin"}       \OtherTok{=}\NormalTok{ beta\_loglin,}
    \StringTok{"R2"}                 \OtherTok{=}\NormalTok{ R2,}
    \StringTok{"p{-}value beta"}       \OtherTok{=}\NormalTok{ Sig}
\NormalTok{  ) }\SpecialCharTok{\%\textgreater{}\%}
  \FunctionTok{mutate}\NormalTok{(}\FunctionTok{across}\NormalTok{(}\FunctionTok{where}\NormalTok{(is.numeric), }\SpecialCharTok{\textasciitilde{}}\FunctionTok{round}\NormalTok{(.x, }\DecValTok{4}\NormalTok{)))}
 
\FunctionTok{cat}\NormalTok{(}\StringTok{"}\SpecialCharTok{\textbackslash{}n}\StringTok{"}\NormalTok{)}
\end{Highlighting}
\end{Shaded}

\begin{Shaded}
\begin{Highlighting}[]
\FunctionTok{print}\NormalTok{(synthese\_finale)}
\end{Highlighting}
\end{Shaded}

\begin{verbatim}
## # A tibble: 3 x 10
##   Zone   `Correlation (r)` `epsilon directe` Interpretation `epsilon croisee TC`
##   <fct>              <dbl>             <dbl> <chr>                         <dbl>
## 1 Urbai~             0.569              1.11 tres elastiqu~                 1.36
## 2 Périu~             0.498              1.06 tres elastiqu~                 1.35
## 3 Rurale             0.413              1.02 tres elastiqu~                 1.33
## # i 5 more variables: `Seuil rupture` <fct>, `Delta epsilon` <dbl>,
## #   `beta log-lin` <dbl>, R2 <dbl>, `p-value beta` <dbl>
\end{verbatim}

\begin{Shaded}
\begin{Highlighting}[]
\CommentTok{\# ══════════════════════════════════════════════════════════════════════════════}
\CommentTok{\#  VISUALISATIONS — SAUVEGARDE}
\CommentTok{\# ══════════════════════════════════════════════════════════════════════════════}
\FunctionTok{cat}\NormalTok{(}\StringTok{"}\SpecialCharTok{\textbackslash{}n}\StringTok{Sauvegarde des graphiques...}\SpecialCharTok{\textbackslash{}n}\StringTok{"}\NormalTok{)}
\end{Highlighting}
\end{Shaded}

\begin{verbatim}
## 
## Sauvegarde des graphiques...
\end{verbatim}

\begin{Shaded}
\begin{Highlighting}[]
\FunctionTok{ggsave}\NormalTok{(}\StringTok{"G1\_introduction\_temporelle.png"}\NormalTok{,  g1, }\AttributeTok{width =} \DecValTok{12}\NormalTok{, }\AttributeTok{height =} \DecValTok{5}\NormalTok{, }\AttributeTok{dpi =} \DecValTok{150}\NormalTok{)}
\FunctionTok{ggsave}\NormalTok{(}\StringTok{"G2\_SQ3\_elasticite\_directe.png"}\NormalTok{,   g2, }\AttributeTok{width =} \DecValTok{10}\NormalTok{, }\AttributeTok{height =} \DecValTok{6}\NormalTok{, }\AttributeTok{dpi =} \DecValTok{150}\NormalTok{)}
\FunctionTok{ggsave}\NormalTok{(}\StringTok{"G3\_SQ1\_seuil\_rupture.png"}\NormalTok{,        g3, }\AttributeTok{width =} \DecValTok{12}\NormalTok{, }\AttributeTok{height =} \DecValTok{6}\NormalTok{, }\AttributeTok{dpi =} \DecValTok{150}\NormalTok{)}
\FunctionTok{ggsave}\NormalTok{(}\StringTok{"G4\_SQ2\_elasticite\_croisee.png"}\NormalTok{,   g4, }\AttributeTok{width =} \DecValTok{10}\NormalTok{, }\AttributeTok{height =} \DecValTok{6}\NormalTok{, }\AttributeTok{dpi =} \DecValTok{150}\NormalTok{)}
\FunctionTok{ggsave}\NormalTok{(}\StringTok{"G5\_SQ3\_regression\_loglin.png"}\NormalTok{,    g5, }\AttributeTok{width =} \DecValTok{12}\NormalTok{, }\AttributeTok{height =} \DecValTok{6}\NormalTok{, }\AttributeTok{dpi =} \DecValTok{150}\NormalTok{)}
\FunctionTok{ggsave}\NormalTok{(}\StringTok{"G6\_SQ2\_report\_modal\_faixas.png"}\NormalTok{,  g6, }\AttributeTok{width =} \DecValTok{12}\NormalTok{, }\AttributeTok{height =} \DecValTok{6}\NormalTok{, }\AttributeTok{dpi =} \DecValTok{150}\NormalTok{)}
 
\CommentTok{\# Grille narrative complète pour slides (ordre narratif : intro {-}\textgreater{} SQ1 {-}\textgreater{} SQ2 {-}\textgreater{} SQ3)}
\NormalTok{grille\_narrative }\OtherTok{\textless{}{-}}\NormalTok{ (g1) }\SpecialCharTok{/}\NormalTok{ (g3 }\SpecialCharTok{|}\NormalTok{ g4) }\SpecialCharTok{/}\NormalTok{ (g2 }\SpecialCharTok{|}\NormalTok{ g5)}
\FunctionTok{ggsave}\NormalTok{(}\StringTok{"SYNTHESE\_DATA\_STORYTELLING.png"}\NormalTok{, grille\_narrative, }\AttributeTok{width =} \DecValTok{20}\NormalTok{, }\AttributeTok{height =} \DecValTok{18}\NormalTok{, }\AttributeTok{dpi =} \DecValTok{150}\NormalTok{)}
 
\FunctionTok{cat}\NormalTok{(}\StringTok{"Graphiques enregistres (7 fichiers).}\SpecialCharTok{\textbackslash{}n}\StringTok{"}\NormalTok{)}
\end{Highlighting}
\end{Shaded}

\begin{verbatim}
## Graphiques enregistres (7 fichiers).
\end{verbatim}

\begin{Shaded}
\begin{Highlighting}[]
\CommentTok{\# ══════════════════════════════════════════════════════════════════════════════}
\CommentTok{\#  EXPORT CSV DÉTAILLÉS (8 fichiers)}
\CommentTok{\# ══════════════════════════════════════════════════════════════════════════════}
\FunctionTok{write.csv}\NormalTok{(synthese\_finale,   }\StringTok{"RESULTATS\_00\_synthese\_finale.csv"}\NormalTok{,         }\AttributeTok{row.names =} \ConstantTok{FALSE}\NormalTok{, }\AttributeTok{fileEncoding =} \StringTok{"UTF{-}8"}\NormalTok{)}
\FunctionTok{write.csv}\NormalTok{(seuil\_table,       }\StringTok{"RESULTATS\_SQ1\_elasticite\_par\_faixa.csv"}\NormalTok{,   }\AttributeTok{row.names =} \ConstantTok{FALSE}\NormalTok{, }\AttributeTok{fileEncoding =} \StringTok{"UTF{-}8"}\NormalTok{)}
\FunctionTok{write.csv}\NormalTok{(seuil\_rupture,     }\StringTok{"RESULTATS\_SQ1\_seuil\_detecte.csv"}\NormalTok{,          }\AttributeTok{row.names =} \ConstantTok{FALSE}\NormalTok{, }\AttributeTok{fileEncoding =} \StringTok{"UTF{-}8"}\NormalTok{)}
\FunctionTok{write.csv}\NormalTok{(cross\_zone,        }\StringTok{"RESULTATS\_SQ2\_elasticite\_croisee.csv"}\NormalTok{,     }\AttributeTok{row.names =} \ConstantTok{FALSE}\NormalTok{, }\AttributeTok{fileEncoding =} \StringTok{"UTF{-}8"}\NormalTok{)}
\FunctionTok{write.csv}\NormalTok{(cross\_faix,        }\StringTok{"RESULTATS\_SQ2\_report\_par\_faixa.csv"}\NormalTok{,       }\AttributeTok{row.names =} \ConstantTok{FALSE}\NormalTok{, }\AttributeTok{fileEncoding =} \StringTok{"UTF{-}8"}\NormalTok{)}
\FunctionTok{write.csv}\NormalTok{(elast\_zone,        }\StringTok{"RESULTATS\_SQ3\_elasticite\_directe.csv"}\NormalTok{,     }\AttributeTok{row.names =} \ConstantTok{FALSE}\NormalTok{, }\AttributeTok{fileEncoding =} \StringTok{"UTF{-}8"}\NormalTok{)}
\FunctionTok{write.csv}\NormalTok{(res\_loglin\_SP95,   }\StringTok{"RESULTATS\_SQ3\_loglin\_SP95\_trafic.csv"}\NormalTok{,     }\AttributeTok{row.names =} \ConstantTok{FALSE}\NormalTok{, }\AttributeTok{fileEncoding =} \StringTok{"UTF{-}8"}\NormalTok{)}
\FunctionTok{write.csv}\NormalTok{(res\_loglin\_TC,     }\StringTok{"RESULTATS\_SQ3\_loglin\_SP95\_TC.csv"}\NormalTok{,         }\AttributeTok{row.names =} \ConstantTok{FALSE}\NormalTok{, }\AttributeTok{fileEncoding =} \StringTok{"UTF{-}8"}\NormalTok{)}
 
\FunctionTok{cat}\NormalTok{(}\StringTok{"CSVs exportes (8 fichiers).}\SpecialCharTok{\textbackslash{}n}\StringTok{"}\NormalTok{)}
\end{Highlighting}
\end{Shaded}

\begin{verbatim}
## CSVs exportes (8 fichiers).
\end{verbatim}

\begin{Shaded}
\begin{Highlighting}[]
\FunctionTok{cat}\NormalTok{(}\StringTok{"}\SpecialCharTok{\textbackslash{}n}\StringTok{SCRIPT EXECUTE AVEC SUCCES !}\SpecialCharTok{\textbackslash{}n\textbackslash{}n}\StringTok{"}\NormalTok{)}
\end{Highlighting}
\end{Shaded}

\begin{verbatim}
## 
## SCRIPT EXECUTE AVEC SUCCES !
\end{verbatim}

\begin{Shaded}
\begin{Highlighting}[]
\CommentTok{\# ══════════════════════════════════════════════════════════════════════════════}
\CommentTok{\#  GUIDE D\textquotesingle{}INTERPRETATION (console) — pour la redaction du rapport}
\CommentTok{\# ══════════════════════════════════════════════════════════════════════════════}
\FunctionTok{cat}\NormalTok{(}\StringTok{"======================================================================}\SpecialCharTok{\textbackslash{}n}\StringTok{"}\NormalTok{)}
\end{Highlighting}
\end{Shaded}

\begin{verbatim}
## ======================================================================
\end{verbatim}

\begin{Shaded}
\begin{Highlighting}[]
\FunctionTok{cat}\NormalTok{(}\StringTok{"  GUIDE D\textquotesingle{}INTERPRETATION POUR LE RAPPORT}\SpecialCharTok{\textbackslash{}n}\StringTok{"}\NormalTok{)}
\end{Highlighting}
\end{Shaded}

\begin{verbatim}
##   GUIDE D'INTERPRETATION POUR LE RAPPORT
\end{verbatim}

\begin{Shaded}
\begin{Highlighting}[]
\FunctionTok{cat}\NormalTok{(}\StringTok{"======================================================================}\SpecialCharTok{\textbackslash{}n}\StringTok{"}\NormalTok{)}
\end{Highlighting}
\end{Shaded}

\begin{verbatim}
## ======================================================================
\end{verbatim}

\begin{Shaded}
\begin{Highlighting}[]
\FunctionTok{cat}\NormalTok{(}\StringTok{"  SQ1 {-} Seuil de rupture : Regarder G3 + tableau SQ1.}\SpecialCharTok{\textbackslash{}n}\StringTok{"}\NormalTok{)}
\end{Highlighting}
\end{Shaded}

\begin{verbatim}
##   SQ1 - Seuil de rupture : Regarder G3 + tableau SQ1.
\end{verbatim}

\begin{Shaded}
\begin{Highlighting}[]
\FunctionTok{cat}\NormalTok{(}\StringTok{"        Si Delta\_max est maximal entre F2 et F3, le seuil}\SpecialCharTok{\textbackslash{}n}\StringTok{"}\NormalTok{)}
\end{Highlighting}
\end{Shaded}

\begin{verbatim}
##         Si Delta_max est maximal entre F2 et F3, le seuil
\end{verbatim}

\begin{Shaded}
\begin{Highlighting}[]
\FunctionTok{cat}\NormalTok{(}\StringTok{"        est empiriquement \textasciitilde{}1.70 EUR (sans l\textquotesingle{}avoir suppose).}\SpecialCharTok{\textbackslash{}n\textbackslash{}n}\StringTok{"}\NormalTok{)}
\end{Highlighting}
\end{Shaded}

\begin{verbatim}
##         est empiriquement ~1.70 EUR (sans l'avoir suppose).
\end{verbatim}

\begin{Shaded}
\begin{Highlighting}[]
\FunctionTok{cat}\NormalTok{(}\StringTok{"  SQ2 {-} Report modal : Regarder G4 + G6.}\SpecialCharTok{\textbackslash{}n}\StringTok{"}\NormalTok{)}
\end{Highlighting}
\end{Shaded}

\begin{verbatim}
##   SQ2 - Report modal : Regarder G4 + G6.
\end{verbatim}

\begin{Shaded}
\begin{Highlighting}[]
\FunctionTok{cat}\NormalTok{(}\StringTok{"        epsilon\_croisee \textgreater{} 0 {-}\textgreater{} TC et voiture sont substituts.}\SpecialCharTok{\textbackslash{}n}\StringTok{"}\NormalTok{)}
\end{Highlighting}
\end{Shaded}

\begin{verbatim}
##         epsilon_croisee > 0 -> TC et voiture sont substituts.
\end{verbatim}

\begin{Shaded}
\begin{Highlighting}[]
\FunctionTok{cat}\NormalTok{(}\StringTok{"        G6 montre si ce report s\textquotesingle{}intensifie au{-}dessus}\SpecialCharTok{\textbackslash{}n}\StringTok{"}\NormalTok{)}
\end{Highlighting}
\end{Shaded}

\begin{verbatim}
##         G6 montre si ce report s'intensifie au-dessus
\end{verbatim}

\begin{Shaded}
\begin{Highlighting}[]
\FunctionTok{cat}\NormalTok{(}\StringTok{"        du seuil detecte en SQ1.}\SpecialCharTok{\textbackslash{}n\textbackslash{}n}\StringTok{"}\NormalTok{)}
\end{Highlighting}
\end{Shaded}

\begin{verbatim}
##         du seuil detecte en SQ1.
\end{verbatim}

\begin{Shaded}
\begin{Highlighting}[]
\FunctionTok{cat}\NormalTok{(}\StringTok{"  SQ3 {-} Facteur geographique : Comparer G2 entre zones.}\SpecialCharTok{\textbackslash{}n}\StringTok{"}\NormalTok{)}
\end{Highlighting}
\end{Shaded}

\begin{verbatim}
##   SQ3 - Facteur geographique : Comparer G2 entre zones.
\end{verbatim}

\begin{Shaded}
\begin{Highlighting}[]
\FunctionTok{cat}\NormalTok{(}\StringTok{"        Rurale {-}\textgreater{} epsilon \textasciitilde{} 0 : contrainte, pas de choix.}\SpecialCharTok{\textbackslash{}n}\StringTok{"}\NormalTok{)}
\end{Highlighting}
\end{Shaded}

\begin{verbatim}
##         Rurale -> epsilon ~ 0 : contrainte, pas de choix.
\end{verbatim}

\begin{Shaded}
\begin{Highlighting}[]
\FunctionTok{cat}\NormalTok{(}\StringTok{"        Urbaine {-}\textgreater{} epsilon plus negatif : alternatives dispo.}\SpecialCharTok{\textbackslash{}n}\StringTok{"}\NormalTok{)}
\end{Highlighting}
\end{Shaded}

\begin{verbatim}
##         Urbaine -> epsilon plus negatif : alternatives dispo.
\end{verbatim}

\begin{Shaded}
\begin{Highlighting}[]
\FunctionTok{cat}\NormalTok{(}\StringTok{"        G5 (log{-}lin) donne le beta structural pour citer}\SpecialCharTok{\textbackslash{}n}\StringTok{"}\NormalTok{)}
\end{Highlighting}
\end{Shaded}

\begin{verbatim}
##         G5 (log-lin) donne le beta structural pour citer
\end{verbatim}

\begin{Shaded}
\begin{Highlighting}[]
\FunctionTok{cat}\NormalTok{(}\StringTok{"        dans le rapport avec p{-}value et R2.}\SpecialCharTok{\textbackslash{}n}\StringTok{"}\NormalTok{)}
\end{Highlighting}
\end{Shaded}

\begin{verbatim}
##         dans le rapport avec p-value et R2.
\end{verbatim}

\begin{Shaded}
\begin{Highlighting}[]
\FunctionTok{cat}\NormalTok{(}\StringTok{"======================================================================}\SpecialCharTok{\textbackslash{}n}\StringTok{"}\NormalTok{)}
\end{Highlighting}
\end{Shaded}

\begin{verbatim}
## ======================================================================
\end{verbatim}

\end{document}
